% AUTO-GENERATED FROM MATH/Math.md (continuation extract) BY build_math_from_md.py.
% Edit Math.md or the builder, then regenerate.
% Options for packages loaded elsewhere
\PassOptionsToPackage{unicode}{hyperref}
\PassOptionsToPackage{hyphens}{url}
\documentclass[
]{article}
\usepackage{xcolor}
\usepackage{amsmath,amssymb}
\setcounter{secnumdepth}{-\maxdimen} % remove section numbering
\usepackage{iftex}
\ifPDFTeX
  \usepackage[T1]{fontenc}
  \usepackage[utf8]{inputenc}
  \usepackage{textcomp} % provide euro and other symbols
\else % if luatex or xetex
  \usepackage{unicode-math} % this also loads fontspec
  \defaultfontfeatures{Scale=MatchLowercase}
  \defaultfontfeatures[\rmfamily]{Ligatures=TeX,Scale=1}
\fi
\usepackage{lmodern}
\ifPDFTeX\else
  % xetex/luatex font selection
\fi
% Use upquote if available, for straight quotes in verbatim environments
\IfFileExists{upquote.sty}{\usepackage{upquote}}{}
\IfFileExists{microtype.sty}{% use microtype if available
  \usepackage[]{microtype}
  \UseMicrotypeSet[protrusion]{basicmath} % disable protrusion for tt fonts
}{}
\makeatletter
\@ifundefined{KOMAClassName}{% if non-KOMA class
  \IfFileExists{parskip.sty}{%
    \usepackage{parskip}
  }{% else
    \setlength{\parindent}{0pt}
    \setlength{\parskip}{6pt plus 2pt minus 1pt}}
}{% if KOMA class
  \KOMAoptions{parskip=half}}
\makeatother
\setlength{\emergencystretch}{3em} % prevent overfull lines
\providecommand{\tightlist}{%
  \setlength{\itemsep}{0pt}\setlength{\parskip}{0pt}}
\usepackage{bookmark}
\IfFileExists{xurl.sty}{\usepackage{xurl}}{} % add URL line breaks if available
\urlstyle{same}
\hypersetup{
  hidelinks,
  pdfcreator={LaTeX via pandoc}}

\author{}
\date{}

\begin{document}

\section{11. Adaptive Selection and Staged
Continuation}\label{adaptive-selection-and-staged-continuation}

\subsection{11.1 Cumulative phase construction for adaptive
selection}\label{cumulative-phase-construction-for-adaptive-selection}

If a shortlisted macro generator \(m\) admits a split family
\(\mathcal C_{\mathrm{split}}(m)=\{m_1,\dots,m_{K_m}\}\), then each
split child inherits the same position \(p\) and therefore defines a
child record \(r_j=(m_j,p)\). For simplicity we drop the \(j\), and
assume each \(m\) below is a unique generator.

Record-first view: the primary object is a candidate-position record
\(r:=(m,p)\), with projections \(\pi_m(m,p):=m\) and \(\pi_p(m,p):=p\).
For each phase \(k \in \{1,2,3\}\), the one-way selector flow is \[
\mathcal R_k\longrightarrow\mathcal S_k\longrightarrow r_k^\star,
\] with \[
\mathcal R_k:=\{(m,p):m\in\mathcal G_k,\ p\in\mathcal P_m^{k}\},
\qquad
\mathcal S_k:=
\{(m,p)\in\mathcal R_k:S_k(m,p)\ge\tau_k\}
\quad\text{or}\quad
\operatorname{Top}_{N_k}\!\left(S_k;\mathcal R_k\right),
\qquad
r_k^\star\in\mathcal S_k.
\] where \(S_k\) denotes the phase-\(k\) score used for shortlist
construction. Generator image and position image:
\(\mathcal G^{(k)}:=\pi_m(\mathcal S_k)\) and
\(\mathcal P_m^{(k)}:=\{p:(m,p)\in\mathcal S_k\}\). Record inheritance:
when later phases inherit the same candidate-position records, the clean
handoff is \(\mathcal R_{k+1}:=\mathcal S_k\). Hence the cumulative
record chain is
\(\mathcal R_3 \subseteq \mathcal R_2 \subseteq \mathcal R_1\), the
cumulative generator chain is
\(\mathcal G^{(3)} \subseteq \mathcal G^{(2)} \subseteq \mathcal G^{(1)} \subseteq \mathcal G\),
the first chain follows from
\(\mathcal R_{k+1}:=\mathcal S_k\subseteq\mathcal R_k\), and the
retained position images inherit analogously via the corresponding
record shortlists \(\mathcal S_1\to\mathcal S_2\to\mathcal S_3\).
Post-shortlist admission, batching, pruning, and beam continuation are
then defined in terms of the final Phase-3 surface.

Seed broadening: in addition, we use a smooth broadening of the initial
candidate record \(\mathcal R\) as expected generator utility saturates,
anti-comensurate to scaffold depth, where the initial seeds are
low-order generators and late-stage seeds are higher-order excitation
and coupling terms.

Chronology: thus the presentation chronologically accords with the
algorithmic implementation: Phase 1 \(\to\) Phase 2 \(\to\) Phase 3,
with each to-arrow filtering its preceding phase.

\subsection{11.2 Phase 1}\label{phase-1}

\subsubsection{11.2.1 Core signal, append position, and one-coordinate
local
model}\label{core-signal-append-position-and-one-coordinate-local-model}

Core record, scaffold, and append slot: \(r=(m,p)\),
\(\mathcal O=(a_1,\dots,a_n)\), \(m\in\mathcal G\), and
\(p\in\mathcal P_m^{1}=\{0,\dots,n\}\). Here \(m\) is a candidate
generator already admitted by the generator gate (so
\(m\in\mathcal G_1\)), and \(p\) is the index of a slot in the
pre-insertion, ordered scaffold \(\mathcal O\). Window bias: note, the
insertion window near the terminal scaffold point is biased to be small
proportional to expected scaffold utility remaining.

Scaffold parameter space and energy map: WLOG consider \(\mathcal O\),
s.t. each \(a\) is assigned one scalar \(\theta (\mod 2 \pi)\), then
define \(\Theta_{\mathcal O}:=\mathbb R^{|\mathcal O|}\),
\(E(\mathcal O,\theta):=\langle\psi(\mathcal O,\theta)|H|\psi(\mathcal O,\theta)\rangle\),
and \(\theta\in\Theta_{\mathcal O}\). Scaffold insertion and parameter
insertion: define
\(\mathcal O\oplus_p m:=(a_1,\dots,a_p,m,a_{p+1},\dots,a_n)\) and
\(\theta\oplus_p\vartheta:=(\theta_1,\dots,\theta_p,\vartheta,\theta_{p+1},\dots,\theta_n)\in\Theta_{\mathcal O\oplus_p m}\).

Write \(|\psi\rangle:=U_{>p}U_{\le p}|\psi_0\rangle\),
\(\widetilde A_r:=U_{>p}A_mU_{>p}^\dagger\),
\(U_r(\alpha):=e^{-i\alpha \widetilde A_r}\), and
\(|\psi_r(\alpha)\rangle:=U_r(\alpha)|\psi\rangle\). Let
\(Q_\psi:=I-|\psi\rangle\langle\psi|\),
\(|d_r\rangle:=\left.\partial_\alpha|\psi_r(\alpha)\rangle\right|_{\alpha=0}=-i\widetilde A_r|\psi\rangle\),
\(\tau_r:=Q_\psi|d_r\rangle\), \(g_r:=2\Re\langle \psi|H|d_r\rangle\),
\(\sigma_r:=\operatorname{SE}\!\left[\widehat g(r)\right]\ge 0\),
\(g_{\mathrm{lcb}}(r):=\max\{|g_r|-z_\alpha\sigma_r,0\}\), and
\(F_{\mathrm{raw}}(r):=\langle \tau_r,\tau_r\rangle+\varepsilon_F\). The
live broad Phase-1 one-coordinate trust-region proxy is \[
\Delta E_{1,\mathrm{TR}}(r)
:=
\max_{|\alpha|\le \rho/\sqrt{F_{\mathrm{raw}}(r)}}
\left[
g_{\mathrm{lcb}}(r)|\alpha|
-\frac12 \lambda_F F_{\mathrm{raw}}(r)\alpha^2
\right].
\] The strict one-coordinate insertion drop is retained as the exact
auxiliary comparator \[
\delta E_1(m,p\mid\mathcal O,\theta)
=
E(\mathcal O,\theta)
-\inf_{\vartheta\in\mathbb R}E\bigl(\mathcal O\oplus_p m,\theta\oplus_p\vartheta\bigr).
\] The infimum notation may be replaced by \(\min\) or
\(\operatorname{Min}\) whenever the minimizer exists. Thus
\(\Delta E_{1,\mathrm{TR}}(r)\) is the live broad-pool second-order
local insertion proxy, whereas \(\delta E_1(m,p\mid\mathcal O,\theta)\)
is the exact one-coordinate insertion drop.
\footnote{User interest: How does window parameterization width scale with cost?}

\subsubsection{11.2.2 Gates, burdens, and feature
ingredients}\label{gates-burdens-and-feature-ingredients}

Base hardware ingredients: each candidate-position pair carries the
proxy triple \(C_{\mathrm{2q}}(m,p)\), \(D(m,p)\), and
\(\Delta K(m,p)\), with fixed nonnegative weights
\(\lambda_{\mathrm{2q}},\lambda_d,\lambda_\theta\).\footnote{Admissibility gate: For each selector stage $k$, let $\Gamma_k:\mathcal G_k^{\mathrm{in}}\to\{0,1\}$ be the generator-level admissibility gate and define $\mathcal G_k^{\mathrm{adm}}:=\{m\in\mathcal G_k^{\mathrm{in}}:\Gamma_k(m)=1\}$. Record domain: the working candidate-record domain is built only from this gated generator pool, $\mathcal R_k:=\{(m,p):m\in\mathcal G_k^{\mathrm{adm}},\ p\in\mathcal P_m^{\mathrm{in},k}\}$. Selector action: all selectors below act on $\mathcal R_k$ itself, so admissibility is enforced by pre-score domain construction rather than post hoc score masking. Tie order: fix once and for all a total order $\prec$ on records (for example, lexicographic order on $(m,p)$). Tie-broken maximizer: for any score $S:\mathcal R_k\to\mathbb R$, define $\arg\max^{\prec}_{r\in\mathcal R_k} S(r):=\min_{\prec}\{r\in\mathcal R_k:S(r)=\max_{u\in\mathcal R_k}S(u)\}$, and similarly for $\arg\min^{\prec}$. Shortlist surface: define $\operatorname{Top}_K^{\prec}(S;\mathcal R_k)$ as the ordered list of the first $K$ records obtained by sorting $\mathcal R_k$ in descending score, with ties broken by $\prec$. Unadorned $\arg\max$, $\arg\min$, and $\operatorname{Top}_K$ below follow this convention.}
Here \(C_{\mathrm{2q}}\) is a two-qubit proxy cost, \(D\) is an expected
depth-induced by \(m\) (at \(p\)) proxy, and \(\Delta K\) is the
marginal parameter count. Positivity assumptions: assume throughout that
\(\lambda_{\mathrm{2q}}\ge 0\), \(C_{\mathrm{2q}}(m,p)\ge 0\), and
\(\lambda_\theta,\lambda_d,\Delta K(m,p),D(m,p)>0\) on admissible
records.
\footnote{Depth and parameter count are constants for our L=2 pareto lean runs.}

\subsubsection{11.2.3 Anticipated utility left and controller
schedule}\label{anticipated-utility-left-and-controller-schedule}

Let \(t\) denote the current selector step, let \(D_{\mathrm{left}}(t)\)
be the remaining controller depth budget, and let
\(\widehat N_{\mathrm{rem}}(t)\) be a heuristic forecast (expected
count) of useful remaining admissions. For gain-process statistics,
define \(\Delta_{\mathrm{adm}}(s):=E(\theta^{(s-1)})-E(\theta^{(s)})\),
\(\Delta_{\mathrm{adm}}^{+}(s):=\operatorname{asinh}\!\left(\frac{\Delta_{\mathrm{adm}}(s)}{\tau_\Delta}\right)\),
and let \(\mathcal I_t^+\) be the index set of past selector steps with
valid admissions used for EWMA/MAD statistics.

Let \(D_{\mathrm{allow}}\) denote the total allowed depth budget, let
\(\ell_t^{\mathrm{plat}}:=\max\{1,\lceil P\,D_{\mathrm{left}}(t)/D_{\mathrm{allow}}\rceil\}\)
and
\(\ell_t^{\mathrm{low}}:=\max\{1,\lceil L\,D_{\mathrm{left}}(t)/D_{\mathrm{allow}}\rceil\}\),
and define the relative-flattening and absolute weak-gain thresholds by
\(\tau_{\mathrm{plat}}^{+}(t):=\bigl(\eta_0+\eta_1(1-D_{\mathrm{left}}(t)/D_{\mathrm{allow}})\bigr)m_t\)
and
\(\tau_{\mathrm{low}}^{+}(t):=\operatorname{asinh}(\tau_{\mathrm{use},t}/\tau_\Delta)\).
The corresponding depth-adaptive streak statistics are
\(\operatorname{plateau\_streak}(t):=\frac{P}{\ell_t^{\mathrm{plat}}}\max\{q\in\{0,\dots,\ell_t^{\mathrm{plat}}\}:\Delta_{\mathrm{adm}}^{+}(t-j+1)\le \tau_{\mathrm{plat}}^{+}(t)\ \forall\,1\le j\le q\}\)
and
\(\operatorname{low\_streak}(t):=\frac{L}{\ell_t^{\mathrm{low}}}\max\{q\in\{0,\dots,\ell_t^{\mathrm{low}}\}:\Delta_{\mathrm{adm}}^{+}(t-j+1)\le \tau_{\mathrm{low}}^{+}(t)\ \forall\,1\le j\le q\}\).

For stagnation, gain, decay, and scale summaries, define
\(u_{\mathrm{stag}}(t):=\operatorname{clip}\!\left(\max\!\left(\frac{\mathrm{plateau\_streak}(t)}{P},\frac{\mathrm{low\_streak}(t)}{L}\right),0,1\right)\),
\(m_t:=\operatorname{EWMA}_{s\in\mathcal I_t^+}\!\bigl(\Delta_{\mathrm{adm}}^{+}(s)\bigr)\),
\(\rho_t:=\operatorname{clip}\!\left(\frac{\operatorname{EWMA}_{\mathrm{recent}}(\Delta_{\mathrm{adm}})}{\operatorname{EWMA}_{\mathrm{older}}(\Delta_{\mathrm{adm}})+\varepsilon},\rho_{\min},1\right)\),
\(\gamma_t:=[-\log\rho_t]_+\), and
\(s_t:=\max\!\left(s_{\min},\operatorname{MAD}_{s\in\mathcal I_t^+}\!\bigl(\Delta_{\mathrm{adm}}^{+}(s)\bigr)\right)\).

For the frontier statistic, let
\(s_{\mathrm{raw},1}^{(1)}(t)\ge s_{\mathrm{raw},2}^{(1)}(t)\ge 0\)
denote the first two order statistics of the current broad Phase-1 score
multiset \(\{S_1(r;t):r\in\mathcal R_1\}\) before any
\(\tau_1(t)\)-thresholding or \(N_1(t)\)-truncation, with the
conventions
\(s_{\mathrm{raw},1}^{(1)}(t)=s_{\mathrm{raw},2}^{(1)}(t)=0\) if
\(\mathcal R_1=\varnothing\) and \(s_{\mathrm{raw},2}^{(1)}(t)=0\) if
\(|\mathcal R_1|=1\). Define \[
u_{\mathrm{front}}(t)
:=
\frac{s_{\mathrm{raw},2}^{(1)}(t)+\varepsilon}{s_{\mathrm{raw},1}^{(1)}(t)+\varepsilon}.
\] Thus \(u_{\mathrm{front}}(t)\in[0,1]\), with larger values
corresponding to a flatter, more crowded current frontier and smaller
values corresponding to a clearly separated leading
record.\footnote{With the empty-pool convention $s_{\mathrm{raw},1}^{(1)}(t)=s_{\mathrm{raw},2}^{(1)}(t)=0$, one gets $u_{\mathrm{front}}(t)=1$. This treats the absence of a usable current frontier as maximal compression rather than as a missing-data null signal. An explicit signal-to-noise controller statistic is a possible implementation here, but it may not be necessary if worsening signal-to-noise already manifests as frontier flattening, which the present expected-utility model scores directly through $u_{\mathrm{front}}(t)$. Here $h_t(j)$ is the conditional probability that usefulness terminates at future offset $j$, given survival through offset $j-1$; accordingly, $S_t(k)$ is the survival probability to the start of offset $k$.}

Then the survival hazard is
\(h_t(j):=\sigma\!\Bigl(\beta_0+\beta_{\mathrm{stag}}u_{\mathrm{stag}}(t)+\beta_{\mathrm{front}}u_{\mathrm{front}}(t)+\eta_h(j-1)\Bigr)\).

\[
S_t(k):=\prod_{j=1}^{k-1}\bigl(1-h_t(j)\bigr),
\qquad
Q_t(k):=\Phi\!\left(
\frac{m_t-\gamma_t(k-1)-\operatorname{asinh}\!\left(\frac{\tau_{\mathrm{use},t}}{\tau_\Delta}\right)}{s_t}
\right).
\] \[
\widehat N_{\mathrm{rem}}(t)
:=
\sum_{k=1}^{D_{\mathrm{left}}(t)}
S_t(k)\,Q_t(k).
\] For pessimistic/optimistic runway envelopes, set
\(h_t^{\mathrm{pess/opt}}(j):=\sigma\!\bigl(\operatorname{logit}h_t(j)\pm\Delta_h\bigr)\),
\(m_t^{\mathrm{pess/opt}}:=m_t\mp\Delta_m\),
\(s_t^{\mathrm{pess/opt}}:=s_t+\Delta_s\) (pess) or
\(s_t^{\mathrm{pess/opt}}:=\max\!\left(s_{\min},s_t-\Delta_s\right)\)
(opt),
\(S_t^{\mathrm{pess/opt}}(k):=\prod_{j=1}^{k-1}\bigl(1-h_t^{\mathrm{pess/opt}}(j)\bigr)\),
and \(Q_t^{\mathrm{pess/opt}}(k):=\Phi\!\left(
\frac{m_t^{\mathrm{pess/opt}}-\gamma_t(k-1)-\operatorname{asinh}\!\left(\frac{\tau_{\mathrm{use},t}}{\tau_\Delta}\right)}{s_t^{\mathrm{pess/opt}}}
\right)\). Then
\(\widehat N_{\mathrm{rem}}^{\mathrm{low/high}}(t):=\sum_{k=1}^{D_{\mathrm{left}}(t)} S_t^{\mathrm{pess/opt}}(k)\,Q_t^{\mathrm{pess/opt}}(k)\),
with
\((\mathrm{low},\mathrm{high})\leftrightarrow(\mathrm{pess},\mathrm{opt})\),
and the confidence ratio is
\(C_t:=\frac{\widehat N_{\mathrm{rem}}^{\mathrm{low}}(t)}{\widehat N_{\mathrm{rem}}^{\mathrm{high}}(t)+\varepsilon}\).

For the controller coordinate and directional transforms, define \[
R_t:=\frac{\widehat N_{\mathrm{rem}}(t)}{D_{\mathrm{left}}(t)+\varepsilon}\in[0,1],
\qquad
e_t:=R_t^{a},
\qquad
g_t:=(1-R_t)^{b},
\qquad
a,b>0.
\] Here \(R_t\) is remaining useful runway as a fraction of the
depth-budget permitted by hardware/run-time constraints, and \(g_t\) is
the late-stage activation coordinate that scales controller variables
depending on age, where age is a normalized quantity relative
to/determined by expected utility of more generators.

Then \(R_t\approx 1\) early and \(R_t\approx 0\) late. Use \(e_t\) for
early-dominant quantities and \(g_t\) for late-dominant quantities. In
particular, define \[
\tau_{\mathrm{pos}}(t):=\tau_{\mathrm{pos}}^{\min}+(\tau_{\mathrm{pos}}^{\max}-\tau_{\mathrm{pos}}^{\min})e_t,
\qquad
\lambda_F(t):=\lambda_F^{\min}+(\lambda_F^{\max}-\lambda_F^{\min})g_t,
\qquad
\nu_{\mathrm{sat}}(t):=1-R_t.
\] Here \(\tau_{\mathrm{pos}}(t)\) is a position-domain parameter: it
controls how broadly the candidate-position fibers \(\mathcal P_m^k(t)\)
or active probe window \(W_{\mathrm{act}}\) may open around append and
active-window positions before phase scores are evaluated. It is not a
score threshold and does not multiply \(S_1\). The separate
inherited-window retention \(\rho_W(t)\) below controls how many old
coordinates are allowed to relax when a surviving record is locally
refit; it changes the geometry used after a record is shortlisted, not
the prior eligibility of a position for scoring.

The quantity \(\widehat N_{\mathrm{rem}}(t)\) is controller telemetry
rather than a physical observable: it summarizes recent selector
progress and remaining useful runway, but it is not itself a hard
admissibility
certificate.\footnote{Interpretation: the model has three layers --- (i) gain process ($m_t,s_t$), (ii) horizon-decay process ($\gamma_t$), and (iii) survival process ($h_t(j)$); together these drive the runway forecast $\widehat N_{\mathrm{rem}}(t)$.}

For inherited refit windows, use the maturity-shrinking retention law \[
\rho_W(t):=\rho_W^{\min}+(\rho_W^{\max}-\rho_W^{\min})e_t,
\qquad
0\le \rho_W^{\min}\le \rho_W^{\max}\le 1,
\] so that the retained old-coordinate fraction is largest early and
contracts as useful runway is exhausted.

The controller doctrine is monotone in scaffold maturity: as
\(R_t\downarrow\) and \(g_t\uparrow\), shot count, prune pressure, and
shortlist pressure increase, while inherited refit breadth, shortlist
breadth, and live phase depth decrease. Phase live/dead decisions are
driven by the runway telemetry \(\widehat N_{\mathrm{rem}}(t)\), not by
raw chronological time. Set \(\Gamma_1^{\mathrm{live}}(t):=1\). For
\(p\in\{2,3\}\), use hysteretic switches \[
\Gamma_p^{\mathrm{live}}(t^+)
:=
\begin{cases}
0,
&
\widehat N_{\mathrm{rem}}^{\mathrm{high}}(t)\le \nu_p^-
\text{ for }h_p\text{ consecutive controller steps},
\\[1mm]
1,
&
\widehat N_{\mathrm{rem}}^{\mathrm{low}}(t)>\nu_p^+,
\\[1mm]
\Gamma_p^{\mathrm{live}}(t^-),
&
\text{otherwise},
\end{cases}
\qquad
0\le \nu_2^-\le \nu_2^+<\nu_3^-\le \nu_3^+.
\] Then enforce
\(\Gamma_3^{\mathrm{live}}(t)\leftarrow\Gamma_2^{\mathrm{live}}(t)\Gamma_3^{\mathrm{live}}(t)\),
so Phase 3 cannot remain live after Phase 2 is nulled. The terminal live
selector phase is \[
k_\star(t):=1+\Gamma_2^{\mathrm{live}}(t)+
\Gamma_2^{\mathrm{live}}(t)\Gamma_3^{\mathrm{live}}(t),
\qquad
r_t^\star\in\arg\max_{r\in\mathcal R_{k_\star}(t)}S_{k_\star}(r;t).
\] Thus early scaffold construction may use the full Phase-1 \(\to\)
Phase-2 \(\to\) Phase-3 chain; when useful runway is nearly exhausted,
Phase 3 is nulled first, then Phase 2 is nulled, leaving the cheaper
greedy Phase-1 selector.

For measurement effort, use \(g_t\) as the canonical maturity-controlled
scheduled-shot baseline. For \(k\in\{1,2,3\}\), define \[
N_{\mathrm{shot},k}^{g}(t)
:=
\left\lceil
N_{\mathrm{shot},k}^{\min}
+
\bigl(N_{\mathrm{shot},k}^{\max}-N_{\mathrm{shot},k}^{\min}\bigr)g_t
\right\rceil,
\qquad
N_{\mathrm{shot},k}^{\min}\le N_{\mathrm{shot},k}^{\max}.
\] Frontier flattening and phase-geometric uncertainty are allowed only
as clipped uplifts on the remaining shot headroom. For Phase 1 set \[
h_1(t)
:=
\operatorname{clip}_{[0,1]}
\!\left(
\beta_{\mathrm{front},1}u_{\mathrm{front}}(t)
\right),
\] while for \(k\in\{2,3\}\), once the phase-\(k\) uncertainty summary
\(u_{\sigma,k}(t)\) is available, set \[
h_k(t)
:=
\operatorname{clip}_{[0,1]}
\!\left(
\beta_{\mathrm{front},k}u_{\mathrm{front}}(t)
+
\beta_{\sigma,k}u_{\sigma,k}(t)
\right),
\qquad
k\in\{2,3\},
\] with nonnegative uplift weights \(\beta_{\mathrm{front},k}\) and
\(\beta_{\sigma,k}\). The bounded scheduled-shot fraction is \[
\Phi_k(t):=g_t+(1-g_t)h_k(t),
\qquad
0\le \Phi_k(t)\le 1,
\] and the scheduled shot law is \[
N_{\mathrm{shot},k}^{\mathrm{sched}}(t):=
\left\lceil
N_{\mathrm{shot},k}^{\min}
+
\bigl(N_{\mathrm{shot},k}^{\max}-N_{\mathrm{shot},k}^{\min}\bigr)\Phi_k(t)
\right\rceil,
\qquad
k\in\{1,2,3\}.
\] Thus
\(h_k(t)=0\Rightarrow N_{\mathrm{shot},k}^{\mathrm{sched}}(t)=N_{\mathrm{shot},k}^{g}(t)\),
\(g_t\le \Phi_k(t)\le 1\), and, holding telemetry fixed, \[
\frac{\partial \Phi_k}{\partial g_t}=1-h_k(t)\ge 0.
\] Therefore \(N_{\mathrm{shot},k}^{g}(t)\) is the doctrinal maturity
floor, while \(u_{\mathrm{front}}(t)\) and \(u_{\sigma,k}(t)\) can only
increase shots above that floor when the selector frontier is flat or
the phase geometry is uncertain. This invariant does not require
irreversible wall-clock monotonicity: if telemetry collapses,
\(N_{\mathrm{shot},k}^{\mathrm{sched}}(t)\) may fall from a
telemetry-inflated value back toward the maturity floor, but it remains
bounded below by \(N_{\mathrm{shot},k}^{g}(t)\).

The actually deployed phase-\(k\) shot count is the live, capped maximum
of this scheduled law and an empirical SNR floor: \[
N_{\mathrm{shot},k}^{\mathrm{snr}}(t)
:=
\left\lceil
\frac{\kappa_k^2\,\widehat\sigma_{t,k}^{\,2}}
{\max\{\widehat\delta_{t,k}^{\,2},\delta_{\mathrm{floor},k}^{\,2}\}}
\right\rceil,
\qquad
N_{\mathrm{shot},k}^{\mathrm{eff}}(t)
:=
\Gamma_k^{\mathrm{live}}(t)
\min\!\left\{
N_{\mathrm{shot},k}^{\mathrm{cap}},
\max\!\left[
N_{\mathrm{shot},k}^{\mathrm{sched}}(t),
N_{\mathrm{shot},k}^{\mathrm{snr}}(t)
\right]
\right\}.
\] Here \(\widehat\delta_{t,k}\) is the useful decision-scale signal,
\(\widehat\sigma_{t,k}\) is its estimator spread, \(\kappa_k\) is the
target SNR, and \(\delta_{\mathrm{floor},k}>0\) prevents pathological
shot explosion when the phase has no resolvable useful gain. If
\(\widehat\delta_{t,k}\) falls below the useful floor, the controller
reads this as no resolvable gain at that phase, not as permission to
spend unbounded shots. The total live measurement burden is \[
N_{\mathrm{tot}}^{\mathrm{eff}}(t)
:=
N_{\mathrm{shot},1}^{\mathrm{eff}}(t)
+
\Gamma_2^{\mathrm{live}}(t)N_{\mathrm{shot},2}^{\mathrm{eff}}(t)
+
\Gamma_2^{\mathrm{live}}(t)\Gamma_3^{\mathrm{live}}(t)N_{\mathrm{shot},3}^{\mathrm{eff}}(t).
\] For phase shortlist thresholds, use the late-stricter monotone law \[
\tau_k^{\mathrm{short}}(t):=\tau_k^{\min}+(\tau_k^{\max}-\tau_k^{\min})g_t,
\qquad
\tau_k^{\min}\le \tau_k^{\max},
\qquad
k\in\{1,2,3\},
\] so that shortlist admission pressure increases as the scaffold
matures.

For phase shortlist breadth, use the live bounded law \[
B_k^{\mathrm{short}}(t):=
\left\lceil
B_k^{\min}+(B_k^{\max}-B_k^{\min})e_t
\right\rceil,
\qquad
B_k^{\min}\le B_k^{\max},
\qquad
k\in\{1,2,3\}.
\] Thus \(R_t\approx1\) early gives low shortlist thresholds and broad
candidate breadth, while \(R_t\approx0\) late gives high shortlist
thresholds and narrow candidate breadth. The invariant is
\(R_t\downarrow\Rightarrow g_t\uparrow\),
\(N_{\mathrm{shot},k}^{\mathrm{sched}}\uparrow\),
\(\tau_k^{\mathrm{short}}\uparrow\), \(B_k^{\mathrm{short}}\downarrow\),
and \(k_\star\downarrow\).

\subsubsection{11.2.4 Phase-local cost and structural burden
definitions}\label{phase-local-cost-and-structural-burden-definitions}

For any current phase-input record family \(\mathcal R_k(t)\), with
\(\mathcal R_1(t):=\mathcal R_1\) and \(k\in\{1,2,3\}\), and any scalar
record feature \(x:\mathcal R_k(t)\to\mathbb R_{\ge 0}\), define the
empirical quantiles
\(q_\alpha^x(t):=\operatorname{Quantile}_\alpha\bigl(\{x(u):u\in\mathcal R_k(t)\}\bigr)\)
and \(m_x(t):=q_{0.5}^x(t)\), together with the robust scales
\(s_x^{\mathrm{MAD}}(t):=1.4826\,\operatorname{median}_{u\in\mathcal R_k(t)}\bigl|x(u)-m_x(t)\bigr|\),
\(s_x^{\mathrm{pct}}(t):=\frac{q_{0.9}^x(t)-q_{0.1}^x(t)}{2.563}\), and
\(s_x(t):=\max\!\bigl(\varepsilon_x,\;s_x^{\mathrm{MAD}}(t),\;s_x^{\mathrm{pct}}(t)\bigr)\).
Thus \(s_x(t)\) is a percentile/MAD hybrid reference scale, with
percentile fallback when the MAD degenerates on sparse discrete
features.

Define the one-sided robust normalized excess by
\(z_x(r;t):=\operatorname{asinh}\!\left(\frac{[x(r)-m_x(t)]_+}{s_x(t)}\right)\),
with \([y]_+:=\max(y,0)\). When a direct ratio form is preferred, use
the reference level \(x_{\mathrm{ref}}(t):=m_x(t)+s_x(t)\).

For a candidate-position record \(r=(m,p)\), let
\(\mathcal L_r:=\operatorname{supp}(m)=\{j:m_j\neq I\}\),
\(\mathrm{wt}(r):=|\mathcal L_r|\), and
\(\#_{X,Y}(r):=\left|\{j\in\mathcal L_r:m_j\in\{X,Y\}\}\right|\).

Define feature-specific robust excesses by
\(z_{\mathrm{wt}}(r;t):=z_x(r;t)\ \text{with }x=\mathrm{wt}\) and
\(z_{XY}(r;t):=z_x(r;t)\ \text{with }x=\#_{X,Y}\). The structural burden
term is \[
D_{\mathrm{struct}}(r;t):=
\beta_{\mathrm{wt}}\,z_{\mathrm{wt}}(r;t)
+\beta_{XY}\,z_{XY}(r;t),
\] with nonnegative coefficients
\(\beta_{\mathrm{wt}},\beta_{XY}\ge 0\).\footnote{Provenance note: generator novelty $G_{\mathrm{new}}(r)$, family novelty $\mathrm{groups}_{\mathrm{new}}(r)$ with single-valued $\mathfrak g:\mathcal G\to\mathcal C$, lifetime/staleness $\mathbf 1_{\mathrm{lifetime}}(r;t)$ with $t_{\mathrm{first}}(m)$ meaning first time ever seen in the run, and the normalized descendants $\bar D_{\mathrm{life}}^{\mathrm{struct}},\bar G_{\mathrm{new}},\bar C_{\mathrm{new}},\bar c,\bar P_{\mathrm{opt}}$ are omitted from the live hardware cost in the present draft and retained only as controller-side provenance/policy bookkeeping. Append-distance pressure is a position-domain construction policy for candidate-position records, not a live hardware-cost term and not a multiplicative Phase-1 score prior.}

For candidate-specific measurement burden, let
\(O_r:=i[H,\widetilde A_r]=\sum_{b\in\mathcal B_r} O_{r,b}\), with
\(O_{r,b}:=\sum_{\ell\in B_{r,b}} c_{r,b\ell} P_{r,b\ell}\), where
\(\mathcal B_r\) is the qubit-wise commuting partition of the
candidate's gradient observable and
\(\mathcal B_r^{\mathrm{new}}\subseteq\mathcal B_r\) is the subset not
already covered by measurement reuse. Here \(\widetilde A_r\) is the
dressed candidate generator at record \(r=(m,p)\), \(\mathcal B_r\) is
the partition of \(O_r=i[H,\widetilde A_r]\) into qubit-wise commuting
Pauli-string measurement groups, and \(\mathcal B_r^{\mathrm{new}}\) is
the subset that still requires new measurement effort. When exact
state-dependent variances are unavailable, use the commutator-based
proxy \[
\widehat v_{r,b}:=\sum_{\ell\in B_{r,b}} |c_{r,b\ell}|^2,
\qquad
\widehat{\bar C}_{\mathrm{shot}}(r):=
\frac{1}{\mathrm{shot}_{\mathrm{ref}}}
\cdot
\frac{
\left(
\sum_{b\in\mathcal B_r^{\mathrm{new}}}\sqrt{\widehat v_{r,b}}
\right)^2
}{
\sigma_{\star}^{2}
}.
\] Here \(\sigma_\star\) is the target standard error for the candidate
gradient estimate. The positive constant
\(\mathrm{shot}_{\mathrm{ref}}\) is a baseline measurement budget
against which candidate-specific shot cost is compared.

For the shared phase-\(k\) cost, define
\(z_{\mathrm{2q}}(r;t):=z_x(r;t)\ \text{with }x=C_{\mathrm{2q}}\),
\(z_d(r;t):=z_x(r;t)\ \text{with }x=D\),
\(z_{\theta}(r;t):=z_x(r;t)\ \text{with }x=\Delta K\),
\(z_{\mathrm{shot}}(r;t):=z_x(r;t)\ \text{with }x=\widehat{\bar C}_{\mathrm{shot}}\),
and \[
K_k(r;t):=
\lambda_{\mathrm{2q}}\,z_{\mathrm{2q}}(r;t)
+\lambda_d\,z_d(r;t)
+\lambda_\theta\,z_{\theta}(r;t)
+\lambda_{\mathrm{shot}}\,z_{\mathrm{shot}}(r;t)
+D_{\mathrm{struct}}(r;t).
\] Explicitly, \[
\begin{aligned}
K_k(r;t)
&=
\lambda_{\mathrm{2q}}\,z_{\mathrm{2q}}(r;t)
+\lambda_d\,z_d(r;t)
+\lambda_\theta\,z_{\theta}(r;t)
+\lambda_{\mathrm{shot}}\,z_{\mathrm{shot}}(r;t)\\
&\qquad
+\beta_{\mathrm{wt}}\,z_{\mathrm{wt}}(r;t)
+\beta_{XY}\,z_{XY}(r;t),
\qquad
k\in\{1,2,3\}.
\end{aligned}
\]

\subsubsection{Algebraic commutation metadata and flat-sector scaffold features}

Before the selector evaluates live energy geometry, the generator pool
may carry static algebraic metadata. This layer identifies support
overlap, bare-generator commutation, and local internally commuting
sectors in the pool or current scaffold. It is a structural prior only:
it is not a dressed runtime commutation claim, and it is not a
substitute for tangent overlap, Hessian coupling, reduced novelty, or
remove-refit recovery.

For bare generator labels define three static metadata entries: the
support-overlap bit, the pairwise commutation bit, and an exactness tag
that records whether the commutation entry was computed exactly or
retained only as approximate metadata: \[
\Omega_{mn}:=\mathbf 1[\operatorname{supp}(m)\cap\operatorname{supp}(n)\neq\varnothing],
\qquad
\mathsf C_{mn}:=\mathbf 1[[A_m,A_n]=0],
\qquad
\mathsf q_{mn}\in\{\mathrm{exact},\mathrm{approx}\}.
\] Here \(\Omega_{mn}\) records whether the two bare generators have
overlapping support, \(\mathsf C_{mn}\) records whether their bare
generators commute, and \(\mathsf q_{mn}\) records the reliability of
that commutation entry: \(\mathrm{exact}\) means the entry was computed
from an exact Pauli expansion or an exact algebraic rule, while
\(\mathrm{approx}\) means only weak metadata was available. The tag
\(\mathsf q_{mn}\) is an exactness label for the pair \((m,n)\), not a
qubit index. For Pauli-word generators
\(A_m=\bigotimes_\ell P_\ell^{(m)}\) and
\(A_n=\bigotimes_\ell P_\ell^{(n)}\), the exact rule is \[
\mathsf C_{mn}
=
\mathbf 1\!\left[
\sum_{\ell\in\operatorname{supp}(m)\cap\operatorname{supp}(n)}
\mathbf 1\!\big[P_\ell^{(m)}P_\ell^{(n)}=-P_\ell^{(n)}P_\ell^{(m)}\big]
\equiv 0\pmod 2
\right],
\qquad
\mathsf q_{mn}=\mathrm{exact}.
\] For non-Pauli composite generators the table is computed from the
expanded Pauli support whenever available; otherwise the corresponding
entries retain \(\mathsf q_{mn}=\mathrm{approx}\) and may be used only
as weak prior evidence.

The static metadata layer keeps three distinct relations: \[
\begin{aligned}
\mathcal E_{\mathrm{disj}}^{\mathrm{comm}}
&:=\{(m,n):\Omega_{mn}=0,\ \mathsf C_{mn}=1\},\qquad
&&\text{support-disjoint commutation},\\
\mathcal E_{\mathrm{flat}}^{\mathrm{comm}}
&:=\{(m,n):\Omega_{mn}=1,\ \mathsf C_{mn}=1\},
&&\text{overlapping flat-sector commutation},\\
\mathcal E_{\mathrm{curv}}^{\mathrm{nc}}
&:=\{(m,n):\Omega_{mn}=1,\ \mathsf C_{mn}=0\},
&&\text{overlapping noncommutation / curvature relation}.
\end{aligned}
\] Disjoint commutation is not redundancy evidence: two support-disjoint
commuting generators may be tangent-orthogonal, geometrically novel, and
simultaneously desirable. The overlapping flat-sector relation is the
relation relevant to low-risk local redundancy and cluster saturation.
The overlapping noncommutation relation is the relation relevant to
frame rotation, curvature-seeking selection, and bridge protection.

Let \(\mathcal N_{\mathrm{loc}}(m)\) be the deterministic locality
neighborhood used for the candidate pool, for example support overlap
plus one-hop hardware or lattice neighbors. This neighborhood is defined
by physical or algebraic locality, not by an interval in the ordered
generator list. Let \(\{\mathfrak C_a\}_a\) denote the internally
commuting clusters returned by a deterministic greedy pass over the
ordered pool, and let \(c(m)=a\) mean that generator \(m\) belongs to
cluster \(\mathfrak C_a\). Equivalently, the pass gives \[
\operatorname{comm\_cluster\_id}(m):=c(m),
\qquad
m,n\in\mathfrak C_a\Longrightarrow \mathsf C_{mn}=1,
\] with tie order inherited from the global record order \(\prec\). The
order \(\prec\) is used only for deterministic ties and reproducible
cluster assignment; it is not the definition of locality.

For Phase-1 stratification, let \(\mathcal C_1(r;t)\subseteq\mathcal G\)
be the local algebraic context of the record \(r=(m,p)\) at selector
step \(t\). This context is built from \(\mathcal N_{\mathrm{loc}}(m)\),
recent scaffold coordinates relevant to the selector, and any inherited
local window information used for the Phase-1 record; it is not an
ordered-list window. Define the exact local relation counts \[
\begin{aligned}
n_{\mathrm{flat}}(r;t)
&:=
\sum_{n\in\mathcal C_1(r;t)}
\mathbf 1[\mathsf q_{mn}=\mathrm{exact}]\,\Omega_{mn}\mathsf C_{mn},\\[1mm]
n_{\mathrm{curv}}(r;t)
&:=
\sum_{n\in\mathcal C_1(r;t)}
\mathbf 1[\mathsf q_{mn}=\mathrm{exact}]\,\Omega_{mn}(1-\mathsf C_{mn}),\\[1mm]
n_{\mathrm{disj}}(r;t)
&:=
\sum_{n\in\mathcal C_1(r;t)}
\mathbf 1[\mathsf q_{mn}=\mathrm{exact}]\,(1-\Omega_{mn})\mathsf C_{mn},\\[1mm]
n_{\mathrm{approx}}(r;t)
&:=
\sum_{n\in\mathcal C_1(r;t)}
\mathbf 1[\mathsf q_{mn}=\mathrm{approx}].
\end{aligned}
\] Let \[
\mathfrak L:=
\{\ell_{\mathrm{flat}},\ell_{\mathrm{curv}},\ell_{\mathrm{disj}},\ell_{\mathrm{mix}}\}
\] be the algebraic lane set, where \(\ell_{\mathrm{flat}}\) denotes
overlapping commuting records, \(\ell_{\mathrm{curv}}\) denotes
overlapping noncommuting records, \(\ell_{\mathrm{disj}}\) denotes
support-disjoint commuting records, and \(\ell_{\mathrm{mix}}\) denotes
mixed or ambiguous records. The support-disjoint lane is a coverage
lane, not a redundancy class. Define the conservative Phase-1 lane map
\[
\ell_1(r;t)
:=
\begin{cases}
\ell_{\mathrm{flat}},
&
n_{\mathrm{flat}}(r;t)>0,\quad
n_{\mathrm{curv}}(r;t)=0,\quad
n_{\mathrm{approx}}(r;t)=0,
\\[1mm]
\ell_{\mathrm{curv}},
&
n_{\mathrm{curv}}(r;t)>0,\quad
n_{\mathrm{flat}}(r;t)=0,\quad
n_{\mathrm{approx}}(r;t)=0,
\\[1mm]
\ell_{\mathrm{disj}},
&
n_{\mathrm{flat}}(r;t)=0,\quad
n_{\mathrm{curv}}(r;t)=0,\quad
n_{\mathrm{disj}}(r;t)>0,\quad
n_{\mathrm{approx}}(r;t)=0,
\\[1mm]
\ell_{\mathrm{mix}},
&
\text{otherwise}.
\end{cases}
\] Thus commutation determines a categorical lane; it does not multiply
the Phase-1 score.

Continuous commutation averages and bridge summaries are not part of the
canonical selector. They are defined only in Appendix 11-S for
telemetry, lane-budget diagnostics, prune-risk bookkeeping, and ablation
studies.

The selection doctrine is local: do not reward noncommutation globally.
Do not penalize consecutive commuting generators merely because they
commute. Use commutation to open Phase-1 lanes, use Phase-2 geometry to
keep or retire inherited lanes for the current selector step, and use
Phase-3 reduced geometry to decide final rank.

Safety convention: do not use commutation as a deletion certificate; do
not treat disjoint commuting generators as redundant; do not use global
commuting degree as a strong downweight; do not allow motif,
commutation, or flat-sector bonuses to rescue weak reduced geometry; and
do not require exact dressed commutation at runtime. Bare commutation is
the precomputed algebraic prior, while tangent overlap, Hessian
coupling, and reduced-window recovery are the effective dressed
geometric diagnostics.

The controller rule is: commutation opens Phase-1 lanes; Phase-2
geometry keeps or retires lanes; Phase-3 geometry decides the final
rank; prune recovery remains categorical through flat-sector versus
curvature-compensated recovery classes.

\subsubsection{11.2.5 Final Phase I Scoring}

For the current scaffold at selector step \(t\), let
\(p_{\mathrm{app}}(t):=|\mathcal O_t|\) be the append slot and define
the normalized append-distance \[
d_{\mathrm{app}}(r;t):=\frac{|p-p_{\mathrm{app}}(t)|}{\max\{1,|\mathcal O_t|\}},
\qquad
0\le d_{\mathrm{app}}(r;t)\le 1.
\] This append-distance is position-domain telemetry. It may be used to
construct or audit the admissible position fibers
\(\mathcal P_m^{1}(t)\), the active probe window
\(W_{\mathrm{act}}(t)\), or the cap \(M_{\mathrm{probe}}\), but it is
not a multiplicative Phase-1 score factor. Position preference is
imposed before scoring by choosing which insertion positions are
eligible for evaluation: \[
\mathcal P_m^{1}(t)\subseteq W_{\mathrm{act}}(t),
\qquad
|\mathcal P_m^{1}(t)|\le M_{\mathrm{probe}},
\qquad
\mathcal R_1(t):=\{(m,p):m\in\mathcal G_1,\ p\in\mathcal P_m^{1}(t)\}.
\] Thus append-like records may receive protected exposure by the
position-window policy, while far internal insertions may be omitted or
delayed by the probe policy. Once a record has entered
\(\mathcal R_1(t)\), the rank score is determined by its geometric gain
and cost burden, not by a separate append-distance multiplier.

For Phase 1, the set \(\mathcal R_1(t)\) is ranked inside algebraic
lanes by the geometric screen, \[
\begin{aligned}
S_{1,\mathrm{geom}}(r;t)
&:=
\frac{\Delta E_{1,\mathrm{TR}}(r)}{1+K_1(r;t)},\\[1mm]
S_{1,\mathrm{rank}}(r;t)
&:=
S_{1,\mathrm{geom}}(r;t),
\qquad r\in\mathcal R_1.
\end{aligned}
\] Here \(\Delta E_{1,\mathrm{TR}}(r)\) is the Phase-1 one-coordinate
trust-region proxy defined in \S 11.2.1, and the phase-1 cost is
evaluated on the current Phase-1 input family. The algebraic layer
assigns the lane \(\ell_1(r;t)\) and controls shortlist coverage; it
does not multiply the Phase-1 rank score. Position-window bias has
already acted by constructing \(\mathcal R_1(t)\); it does not reweight
records after their Phase-1 quantities have been computed. In
particular, Phase 1 does not penalize a candidate merely because it
commutes with the last generator, and support-disjoint commutation is
not a redundancy signal. This Phase-1 rank score is intentionally broad
as a screening surface: it privileges cheap local second-order insertion
estimates while the controller still has long useful runway. The
corresponding broad Phase-1 probing/evaluation budget is set by the
baseline controller shot law \(N_{\mathrm{shot},1}(t)\).

In the present draft, the live Phase-1 numerator is the one-coordinate
trust-region proxy \(\Delta E_{1,\mathrm{TR}}(r)\), so that the broad
screen remains in the same gradient-plus-second-order family later
refined in Phases 2 and 3. The strict one-coordinate insertion drop
\(\delta E_1(r\mid\mathcal O,\theta)\) is retained only as an auxiliary
exact comparator for selective frontier confirmation, ablation studies,
or future confirm-stage refinement. A broader early-stage local-refit
numerator remains only a non-live hook. To make such a refinement live,
one would first have to specify (i) the admissible refit-window family,
(ii) the corresponding constrained local-refit drop, and (iii) a
controller activation law independent of the current raw Phase-1 score
multiset, so that the live Phase-1 law is not self-referential.

\subsubsection{11.2.6 Position probing, lane-wise trough detection, and
effective
selector}\label{position-probing-lane-wise-trough-detection-and-effective-selector}

The Phase-1 scoring mechanism produces a stratified record shortlist for
further refined scoring. For each algebraic lane \(\ell\in\mathfrak L\),
define \[
\mathcal R_{1,\ell}(t):=\{r\in\mathcal R_1(t):\ell_1(r;t)=\ell\},
\qquad
\mathcal A_{1,\ell}(t):=\{r\in\mathcal R_{1,\ell}(t):S_{1,\mathrm{rank}}(r;t)\ge \tau_1(t)\}.
\] Let \(B_{1,\ell}(t)\in\mathbb N\) be the Phase-1 shortlist budget
assigned to lane \(\ell\). The default lane-budget doctrine is
asymmetric. The mixed lane receives protected exposure because mixed can
mean unresolved or approximate algebraic evidence rather than weak
evidence. For \(\mathcal R_{1,\mathrm{mix}}(t)\neq\varnothing\), define
\[
\bar f_{\mathrm{approx}}^{\mathrm{mix}}(t)
:=
\frac{
\sum_{r\in\mathcal R_{1,\mathrm{mix}}(t)} n_{\mathrm{approx}}(r;t)
}{
\sum_{r\in\mathcal R_{1,\mathrm{mix}}(t)}
\left[n_{\mathrm{flat}}(r;t)+n_{\mathrm{curv}}(r;t)+n_{\mathrm{disj}}(r;t)+n_{\mathrm{approx}}(r;t)\right]
\varepsilon
}.
\] A canonical protected mixed budget is \[
B_{1,\mathrm{mix}}(t)
:=
\min\!\left\{
B_{1,\mathrm{mix}}^{\max},
B_{1,\mathrm{mix}}^{\min}
+
\left\lceil
\kappa_{\mathrm{mix}}\,
\bar f_{\mathrm{approx}}^{\mathrm{mix}}(t)\,
B_1^{\mathrm{short}}(t)
\right\rceil
\right\},
\qquad
B_{1,\mathrm{mix}}^{\min}\ge 1.
\] Thus approximate metadata can increase mixed-lane exposure, but it
does not increase any record's score. The disjoint lane has the
complementary default: it receives modest Phase-1 exposure because
support-disjoint commutation is not redundancy evidence, but it is often
less locally corrective than overlapping flat or curvature evidence. A
typical default satisfies \[
B_{1,\mathrm{disj}}^{\min}
\le
B_{1,\mathrm{flat}}^{\min},
\qquad
B_{1,\mathrm{disj}}^{\min}
\le
B_{1,\mathrm{curv}}^{\min},
\] while still keeping \(B_{1,\mathrm{disj}}^{\min}>0\) whenever
disjoint exploration is enabled. These lane budgets control exposure
only; records inside every lane are still ordered by
\(S_{1,\mathrm{rank}}\).

Order the records in \(\mathcal A_{1,\ell}(t)\) as \[
r_{\ell,(1)}^{(1)},\dots,r_{\ell,(M_\ell)}^{(1)},
\qquad
s_{\ell,j}^{(1)}(t):=S_{1,\mathrm{rank}}\!\left(r_{\ell,(j)}^{(1)};t\right),
\] with
\(s_{\ell,1}^{(1)}(t)\ge s_{\ell,2}^{(1)}(t)\ge\cdots\ge s_{\ell,M_\ell}^{(1)}(t)\),
ties resolved by \(\prec\). For \(M_\ell\ge2\), define the lane-local
frontier ratio \[
u_{\mathrm{front},k}^{(1,\ell)}(t)
:=
\frac{s_{\ell,k+1}^{(1)}(t)+\varepsilon}{s_{\ell,k}^{(1)}(t)+\varepsilon},
\qquad
k=1,\dots,M_\ell-1.
\] With \(0<\eta_{1,\ell}<1\), set \[
\begin{aligned}
k_{1,\ell}^\star(t)&:=
\begin{cases}
0, & M_\ell=0,\\[4pt]
\min\!\left\{k\in\{1,\dots,\min(M_\ell-1,B_{1,\ell}(t)-1)\}:u_{\mathrm{front},k}^{(1,\ell)}(t)\le \eta_{1,\ell}\right\},
& \text{if this set is nonempty},\\[6pt]
\min\{M_\ell,B_{1,\ell}(t)\}, & \text{otherwise},
\end{cases}\\
\mathcal S_{1,\ell}(t)&:=\{r_{\ell,(j)}^{(1)}:1\le j\le k_{1,\ell}^\star(t)\}.
\end{aligned}
\] The retained Phase-1 shortlist is the lane union \[
\mathcal S_1(t):=\bigcup_{\ell\in\mathfrak L}\mathcal S_{1,\ell}(t).
\] The retained model set is
\(\mathcal G^{(1)}:=\pi_m(\mathcal S_1(t))\), the retained position
fibers are
\(\mathcal P_m^{(1)}:=\{p:(m,p)\in\mathcal S_1(t)\}\ \forall m\in\mathcal G^{(1)}\),
and the Phase-2 record pool is \(\mathcal R_2(t):=\mathcal S_1(t)\).
Thus Phase 1 gives each algebraic lane controlled exposure, while the
ranking within each lane remains geometric-cost based.

\subsection{11.3 Phase II}\label{phase-ii}

Phase II is the raw geometric selector stage. It acts on the Phase-1
survivors, restores the coarse geometric signal that distinguishes
directional usefulness beyond insertion-drop alone, and produces the
filtered shortlist \(\mathcal S_2\).

\subsubsection{11.3.1 Raw geometry, restored novelty, and auxiliary
overlap
window}\label{raw-geometry-restored-novelty-and-auxiliary-overlap-window}

The Phase-2 selector works over candidate-position records inherited
from Phase 1. In the record-inheritance mode used here, the inherited
and admissible domain is \(\mathcal R_2(t):=\mathcal S_1(t)\). Each
inherited record carries its Phase-1 lane label \[
\ell_2(r;t):=\ell_1(r;t),
\qquad
\mathcal R_{2,\ell}(t):=\{r\in\mathcal R_2(t):\ell_2(r;t)=\ell\},
\qquad
\ell\in\mathfrak L.
\] For \(r=(m,p)\), let \(W(r;t)\subseteq\{1,\dots,n+1\}\) denote the
inherited overlap window attached to \(r\). In the present three-phase
architecture this window enters the Phase-2 selector through the Fubini
tangent span of the current horizontal tangent family; the same
inherited window is then carried forward to Phase 3, where
reduced-window elimination and final reranking are
performed.\footnote{Here $p$ indexes the cut after $U_p$. Thus $U_{>p}:=U_K\cdots U_{p+1}$ and $U_{\le p}:=U_p\cdots U_1$, and insertion at cut $p$ means $ |\psi_r(\alpha)\rangle = U_{>p}e^{-i\alpha A_m}U_{\le p}|\psi_0\rangle = U_K\cdots U_{p+1}\,e^{-i\alpha A_m}\,U_p\cdots U_1|\psi_0\rangle $. So the new generator is inserted between $U_p$ and $U_{p+1}$.}

Reusing the record-local insertion objects already introduced in Phase
1, let \(|d_r^{(2)}\rangle:=-\widetilde A_r^2|\psi\rangle\). The
gradient estimator uncertainty is already accounted for through the
lower-confidence gradient
\(g_{\mathrm{lcb}}(r)=\max\{|g_r|-z_\alpha\sigma_r,0\}\), so the live
Phase-2 and Phase-3 scores do not include an additional multiplicative
confidence
factor.\footnote{Let $\widehat g(r)$ denote the empirical gradient estimator for record $r=(m,p)$ obtained from $N_r$ independent shot-level repetitions. If $X^{(r)}_1,\dots,X^{(r)}_{N_r}$ are the repetition-level gradient samples, then $\widehat g(r):=\frac{1}{N_r}\sum_{i=1}^{N_r} X^{(r)}_i$ and $\sigma_r:=\operatorname{SE}\!\big[\widehat g(r)\big]=\sqrt{\operatorname{Var}\!\big(\widehat g(r)\big)}$. Under independent repetitions, $\operatorname{Var}\!\big(\widehat g(r)\big)=\frac{\operatorname{Var}\!\big(X^{(r)}_1\big)}{N_r}$, so in practice one estimates $\sigma_r\approx \frac{s_r}{\sqrt{N_r}}$, where $s_r$ is the sample standard deviation of the repetition-level gradients. Thus $\sigma_r$ is the standard error of the mean gradient estimate, not the sample standard deviation itself.}
For phase-dependent shot allocation in later phases, define \[
u_{\sigma,k}(t):=
\operatorname{clip}\!\left(
\operatorname{median}_{r\in\mathcal R_k(t)}
\frac{z_\alpha\,\sigma_r}{|g_r|+\varepsilon_g},
\,0,\,1
\right),
\qquad
k\in\{2,3\},
\] where in Phase 2 the record-local gradient and uncertainty estimates
are evaluated using \(N_{\mathrm{shot},2}(t)\), and in Phase 3 the pair
\((g_r,\sigma_r)\) is inherited from the Phase-II raw geometry carried
by the surviving records.

The raw Phase-2 tangent-span quantities are \[
\begin{aligned}
F_{\mathrm{raw}}(r)&:=\langle \tau_r,\tau_r\rangle+\varepsilon_F,
\qquad \varepsilon_F>0,\\
(Q_r^{(2)})_{ij}
&:=
\Re\langle \tau_i,\tau_j\rangle,
\qquad i,j\in W(r;t),\\
(q_r^{(2)})_i
&:=
\Re\langle \tau_i,\tau_r\rangle,
\qquad i\in W(r;t),\\
z_r^{(2)}
&:=
\frac{(q_r^{(2)})^\top
\bigl(Q_r^{(2)}+\lambda_{2,\mathrm{nov}}I\bigr)^{-1}
q_r^{(2)}}
{F_{\mathrm{raw}}(r)},
\qquad \lambda_{2,\mathrm{nov}}>0,\\
\mathcal N_2(r;t)
&:=\operatorname{clip}_{[0,1]}\!\left(1-z_r^{(2)}\right),\\
h_r&:=2\Re\!\left(\langle d_r|H|d_r\rangle+\langle \psi|H|d_r^{(2)}\rangle\right),\\
\Delta E_{\mathrm{TR}}^{\mathrm{raw}}(r)
&:=\max_{|\alpha|\le \rho/\sqrt{F_{\mathrm{raw}}(r)}}
\left[g_{\mathrm{lcb}}(r)|\alpha|-\frac12\max\bigl(h_r,0\bigr)\alpha^2\right].
\end{aligned}
\] Here \(F_{\mathrm{raw}}(r)\) is the raw candidate tangent norm,
\(Q_r^{(2)}\) is the inherited-window Fubini Gram matrix, \(q_r^{(2)}\)
is the candidate-to-window tangent covariance vector,
\(\mathcal N_2(r;t)\) is the raw Phase-2 span-novelty score, \(h_r\) is
the raw local curvature, and
\(\Delta E_{\mathrm{TR}}^{\mathrm{raw}}(r)\) is the raw trust-region
gain before reduced-window
elimination.\footnote{The archaic Phase-II root used the pairwise proxy $1-\max_{j\in W(r;t)}|\langle\tau_j,\tau_r\rangle|/(\|\tau_j\|\,\|\tau_r\|)$. That proxy is retained only as a cheap diagnostic or fallback in implementations that do not materialize $Q_r^{(2)}$. It is not the canonical manuscript novelty factor, because pairwise low overlap can miss collective span redundancy.}

The authoritative Phase-II score is \[
S_2(r;t)
:=
\frac{\Delta E_{\mathrm{TR}}^{\mathrm{raw}}(r)\,\mathcal N_2(r;t)}{1+K_2(r;t)},
\qquad r\in\mathcal R_2.
\] Thus Phase II remains the raw geometric selector stage: it restores
directional discrimination beyond insertion-drop alone, uses the
lower-confidence trust-region gain to handle gradient noise, tests
candidate independence against the inherited tangent span, and filters
the Phase-1 survivors before any reduced-window refit quantity is made
authoritative. Commutation no longer modifies the Phase-2 score; it only
supplies the inherited lane label whose current usefulness is tested by
Phase-2 geometry.

For continuity with the older local-refit notation, retain the auxiliary
windowed drop \[
\mathcal C_W(r;\theta):=\{\widetilde\theta\in\Theta_{\mathcal O\oplus_p m}:\widetilde\theta_j=(\theta\oplus_p 0)_j\ \forall j\notin W(r;t)\}
\] and \[
\delta E_2(r\mid\mathcal O,\theta):=
E(\mathcal O,\theta)-\inf_{\widetilde\theta\in\mathcal C_W(r;\theta)}E(\mathcal O\oplus_p m,\widetilde\theta).
\] This \(\delta E_2\) quantity is retained only as an auxiliary
comparator for the inherited window. It may be used for notation
continuity, diagnostics, or ablation studies, but it does not define a
separate selector phase, it does not enter \(S_2(r;t)\), and it does not
modify the Phase-II shortlist rule. Any authoritative reduced-window
rerank remains deferred to Phase 3.

\subsubsection{11.3.2 Inherited lane survival and shortlist
rule}\label{inherited-lane-survival-and-shortlist-rule}

Phase 2 first ranks records inside each inherited lane, then ranks lanes
against one another through a geometry-only lane-health statistic. For
each \(\ell\in\mathfrak L\), let \[
M_{2,\ell}(t):=|\mathcal R_{2,\ell}(t)|.
\] If \(M_{2,\ell}(t)>0\), order \(\mathcal R_{2,\ell}(t)\) as \[
r_{\ell,(1)}^{(2)},\dots,r_{\ell,(M_{2,\ell})}^{(2)},
\qquad
s_{\ell,j}^{(2)}(t):=S_2\!\left(r_{\ell,(j)}^{(2)};t\right),
\] with
\(s_{\ell,1}^{(2)}(t)\ge s_{\ell,2}^{(2)}(t)\ge\cdots\ge s_{\ell,M_{2,\ell}}^{(2)}(t)\),
ties resolved by \(\prec\). The default lane-health statistic is the
best resolved raw-geometric score in the lane, \[
H_{2,\ell}(t):=
\begin{cases}
s_{\ell,1}^{(2)}(t),&M_{2,\ell}(t)>0,\\
0,&M_{2,\ell}(t)=0.
\end{cases}
\] This default encodes the rule that one strong Phase-2 candidate is
enough to keep a lane alive for the current selector step. A top-few
mean may be used as a robustness ablation, but it is not the canonical
lane-health statistic. Define the relative lane health \[
H_{2,\star}(t):=\max_{\ell'\in\mathfrak L}H_{2,\ell'}(t),
\qquad
\widehat H_{2,\ell}(t):=
\begin{cases}
H_{2,\ell}(t)/H_{2,\star}(t),&H_{2,\star}(t)>0,\\
0,&H_{2,\star}(t)=0.
\end{cases}
\] Let \(\tau_2^{\mathrm{abs}}(t)\ge0\) be the common absolute
lane-survival scale and \(\tau_2^{\mathrm{rel}}(t)\in[0,1]\) the common
relative lane-survival scale. Lane-specific patience factors
\(c_\ell^{\mathrm{abs}}(t)>0\) and \(c_\ell^{\mathrm{rel}}(t)>0\) define
\[
\tau_{2,\ell}^{\mathrm{abs}}(t):=c_\ell^{\mathrm{abs}}(t)\tau_2^{\mathrm{abs}}(t),
\qquad
\tau_{2,\ell}^{\mathrm{rel}}(t):=
\operatorname{clip}_{[0,1]}\!\left(c_\ell^{\mathrm{rel}}(t)\tau_2^{\mathrm{rel}}(t)\right).
\] Smaller \(c_\ell^{\mathrm{abs}}\) or \(c_\ell^{\mathrm{rel}}\) gives
lane \(\ell\) more patience; larger values require stronger geometric
evidence. These factors may be fixed, scheduled by scaffold maturity, or
learned by the outer policy, but they affect only lane survival and lane
budgets, not the intra-lane record ranking.

The live-lane gate is \[
\Gamma_{2,\ell}^{\mathrm{live}}(t)
:=
\mathbf 1[M_{2,\ell}(t)>0]\,
\mathbf 1[H_{2,\ell}(t)\ge\tau_{2,\ell}^{\mathrm{abs}}(t)]\,
\mathbf 1[\widehat H_{2,\ell}(t)\ge\tau_{2,\ell}^{\mathrm{rel}}(t)].
\] Lane failure is local to the current selector step:
\(\Gamma_{2,\ell}^{\mathrm{live}}(t)=0\) does not delete the lane from
the generator universe and does not certify algebraic redundancy.

For live lanes, define the lane-local record acceptance set by a record
floor and an intra-lane closeness rule, \[
\mathcal A_{2,\ell}(t)
:=
\left\{
r\in\mathcal R_{2,\ell}(t):
S_2(r;t)\ge\tau_2(t),\
S_2(r;t)\ge \eta_{2,\ell}^{\mathrm{in}}(t)s_{\ell,1}^{(2)}(t)
\right\},
\] where \(0\le\eta_{2,\ell}^{\mathrm{in}}(t)\le1\). Let
\(B_{2,\ell}(t)\in\mathbb N\) be the Phase-2 budget assigned to lane
\(\ell\). The lane-wise Phase-2 shortlist is \[
\mathcal S_{2,\ell}(t)
:=
\begin{cases}
\operatorname{Top}_{B_{2,\ell}(t)}\!\left(S_2;\mathcal A_{2,\ell}(t)\right),
&
\Gamma_{2,\ell}^{\mathrm{live}}(t)=1,\\[1mm]
\varnothing,
&
\Gamma_{2,\ell}^{\mathrm{live}}(t)=0.
\end{cases}
\] The retained Phase-2 shortlist is \[
\mathcal S_2(t):=\bigcup_{\ell\in\mathfrak L}\mathcal S_{2,\ell}(t).
\] Of course, \(\mathcal G^{(2)}:=\pi_m(\mathcal S_2(t))\), and
\(\mathcal P_m^{(2)}:=\{p:(m,p)\in\mathcal S_2(t)\}\). Thus Phase 1
gives algebraic lanes a chance, Phase 2 keeps or retires inherited lanes
using raw geometry, and Phase 3 receives only the surviving records.

\subsection{11.4 Phase 3 reduced-window rerank and post-shortlist
admission}\label{phase-3-reduced-window-rerank-and-post-shortlist-admission}

This section refines the surviving Phase-2 records with the more
expensive reduced-window geometry, then uses that final rerank surface
for shortlist formation, batched admission, pruning, and continuation.

\subsubsection{11.4.1 Reduced-window geometry, curvature, and
trust-region
gain}\label{reduced-window-geometry-curvature-and-trust-region-gain}

Phase 3 starts from the Phase-II raw quantities \(\tau_r\), \(g_r\),
\(h_r\), \(g_{\mathrm{lcb}}(r)\), \(F_{\mathrm{raw}}(r)\), and
\(\mathcal N_2(r;t)\), then refines them using the inherited refit
window. The raw span-novelty score \(\mathcal N_2(r;t)\) is not
recomputed here; instead the candidate direction is reduced against the
inherited window before the expensive rerank surface is formed.

\subsubsection{11.4.2 Refit window, reduced geometry, and Phase-3 record
family}\label{refit-window-reduced-geometry-and-phase-3-record-family}

In the record-inheritance mode used here, the Phase-3 input set is
\(\mathcal R_3(t):=\mathcal S_2(t)\). The inherited refit window
attached to position \(p\) at selector step \(t\) is denoted \(W(p;t)\),
with decomposition
\(W(p;t)=W_{\mathrm{new}}^{(p)}\cup W_{\mathrm{age}}^{(p)}\), where
\(W_{\mathrm{new}}^{(p)}\) keeps the newest coordinates and
\(W_{\mathrm{age}}^{(p)}\) preserves the maturity-scaled set of oldest
retained old coordinates defined later by the branch-local admission
order. For each shortlisted record \(r=(m,p)\), let \(W_r:=W(r;t)\) be
the inherited window that would actually be refit if \(m\) were admitted
at position \(p\), and let \(\{\tau_j\}_{j\in W_r}\) be the current
horizontal tangents in that inherited window. Phase 3 evaluates its
inherited reduced-window geometry under the later-phase controller shot
law \(N_{\mathrm{shot},3}(t)\).

Using the Phase-II raw tangent data, define
\((Q_r)_{jk}=\Re\langle \tau_j,\tau_k\rangle\),
\((q_r)_j=\Re\langle \tau_j,\tau_r\rangle\),
\(E_0:=E(\theta;\mathcal O)\),
\((b_r)_j:=\left.\partial_\alpha\partial_{\theta_j}E(\alpha,\delta\theta_{W_r};r)\right|_{\alpha=0,\,\delta\theta_{W_r}=0}\)
for \(j\in W_r\), and
\((H_{W_r})_{jk}:=\left.\partial_{\theta_j}\partial_{\theta_k}E(\alpha,\delta\theta_{W_r};r)\right|_{\alpha=0,\,\delta\theta_{W_r}=0}\)
for \(j,k\in W_r\). The local reduced-window energy model is \[
E(\alpha,\delta\theta_{W_r};r)
\approx
E_0+g_r\alpha+\tfrac12 h_r\alpha^2+\alpha b_r^\top\delta\theta_{W_r}+\tfrac12\delta\theta_{W_r}^\top H_{W_r}\delta\theta_{W_r}.
\] For reduction and reduced novelty,
\(\mathbf R_r=\tfrac12(H_{W_r}+H_{W_r}^{\top})+\lambda_H I\),
\(\widetilde h_r=h_r-b_r^\top \mathbf R_r^{-1}b_r\),
\(F_r^{\mathrm{red}}=F_{\mathrm{raw}}(r)-2q_r^\top \mathbf R_r^{-1}b_r+b_r^\top \mathbf R_r^{-1}Q_r\mathbf R_r^{-1}b_r\),
\(F_{r,\mathrm{safe}}^{\mathrm{red}}=\max\!\left(F_r^{\mathrm{red}},\varepsilon_{\mathrm{red}}\right)\)
with \(\varepsilon_{\mathrm{red}}>0\),
\(q_r^{\mathrm{red}}=q_r-Q_r\mathbf R_r^{-1}b_r\), and
\(z_r:=\frac{(q_r^{\mathrm{red}})^\top (Q_r+\varepsilon_{\mathrm{nov}}I)^{-1} q_r^{\mathrm{red}}}{F_{r,\mathrm{safe}}^{\mathrm{red}}}\),
\(\mathcal N_3(r;t):=\operatorname{clip}_{[0,1]}\!\left(1-z_r\right)\).
The reduced trust-region gain is \[
\Delta E_{\mathrm{TR}}(r)
:=
\max_{|\alpha|\le \rho/\sqrt{F_{r,\mathrm{safe}}^{\mathrm{red}}}}
\left[g_{\mathrm{lcb}}(r)|\alpha|-\tfrac12\max(\widetilde h_r,0)\alpha^2\right].
\] This reduced tangent is an energy-relaxed tangent, not a pure
orthogonal projection: the Hessian block \(\mathbf R_r\) encodes the
local inherited-window relaxation used to ask how the candidate would
look after nearby old coordinates are allowed to
adjust.\footnote{Equivalently, the reduction mixes tangent-space geometry with the local energetic relaxation model. That is the intended Phase-III semantics; pure span novelty is only an approximation to this energy-relaxed geometry.}

\subsubsection{11.4.3 Phase-3 rerank score and effective selector}

The authoritative Phase-3 rerank score is the reduced-window geometric
score alone: \[
S_3(r;t)
:=
\frac{\Delta E_{\mathrm{TR}}(r)\,\mathcal N_3(r;t)}{1+K_3(r;t)},
\qquad r\in\mathcal R_3(t).
\] Thus Phase 3 is geometry-authoritative: the reduced-window
lower-confidence trust-region gain, reduced novelty, and phase-local
burden determine the final rank. The inherited lane label
\(\ell_3(r;t):=\ell_2(r;t)\), commutation telemetry, operator-family
labels, and motif descriptors may be logged, may inform future
lane-budget policy, or may resolve deterministic ties after geometric
equality. They do not enter the canonical Phase-3 score as additive or
multiplicative rescue terms.

The retained Phase-3 shortlist is obtained from
\(\mathcal A_3(t):=\{r=(m,p)\in\mathcal R_3(t):S_3(r;t)\ge \tau_3(t)\}=\{r_{(1)},\dots,r_{(M)}\}\),
with \(M:=|\mathcal A_3(t)|\). Let \(s_k^{(3)}(t):=S_3(r_{(k)};t)\),
such that \(s^{(3)}_1(t)\ge s^{(3)}_2(t)\ge\cdots\ge s^{(3)}_M(t)\).
Rate the candidate's score distinction by
\(u_{\mathrm{front},k}^{(3)}(t):=\frac{s^{(3)}_{k+1}(t)+\varepsilon}{s^{(3)}_{k}(t)+\varepsilon}\),
with \(k=1,\dots,M-1\) and \(0<\eta_3<1\). Then \[
\begin{aligned}
k_3^\star(t)&:=
\begin{cases}
0, & M=0,\\[4pt]
\min\!\left\{k\in\{1,\dots,\min(M-1,N_3(t)-1)\}:u_{\mathrm{front},k}^{(3)}(t)\le \eta_3\right\},
& \text{if this set is nonempty},\\[6pt]
\min\{M,N_3(t)\}, & \text{otherwise},
\end{cases}\\
\mathcal D_3(t)&:=\{r_{(j)}:1\le j\le k_3^\star(t)\},
\qquad
\mathcal S_3(t):=\mathcal A_3(t)\cap \mathcal D_3(t)=\mathcal D_3(t).
\end{aligned}
\] Of course, \(\mathcal G^{(3)}:=\pi_m(\mathcal S_3(t))\), and
\(\mathcal P_m^{(3)}:=\{p:(m,p)\in\mathcal S_3(t)\}\).

\subsubsection{11.4.4 Reduced-plane geometric batch selection after
shortlisting}\label{reduced-plane-geometric-batch-selection-after-shortlisting}

After the Phase-3 shortlisting has produced the record family
\(\mathcal S_3(t)\subseteq \mathcal R_3(t)\), batching is resolved by
reduced-state-space geometry on a common inherited tangent plane, rather
than by heuristic pairwise compatibility penalties. Support-overlap,
commutation, and related proxy features may still be used as optional
cheap prescreens or route-level compatibility filters, but they are not
part of the principled batching law below unless the chosen route
explicitly declares them to be hard cross-record constraints. Assume
throughout that \(\mathcal S_3(t)\neq\varnothing\); otherwise no Phase-3
batch is formed.

Disjoint support plays a different role in batching than it plays in
Phase 1. The relation \(\Omega_{mn}=0\) is not a Phase-1 redundancy
signal and does not by itself justify a large exploratory shortlist
budget. However, after two records have survived Phase 2 and Phase 3
geometry, support-disjoint structure is favorable batching evidence: it
suggests lower structural conflict, weaker direct measurement overlap,
and a better chance that the two reduced directions can be co-admitted
without destructive local interference. Thus disjoint records receive
modest protected exposure before geometry, but may become preferred
low-conflict batch partners after geometry has validated them.

Near-degenerate shell and structural feasibility: let
\(s_t(r):=S_3(r;t)\) and
\(r_{\max}(t)\in \arg\max_{r\in \mathcal S_3(t)} s_t(r)\). Fix a
near-degenerate shell parameter \(\eta_{\mathrm{nd}}\in(0,1]\) and a
batch cap \(B_{\max}(t)\in\mathbb N\). Define the near-degenerate record
shell
\(\mathcal N_3(t):=\left\{ r\in \mathcal S_3(t): s_t(r)\ge \eta_{\mathrm{nd}}\,s_t\!\bigl(r_{\max}(t)\bigr) \right\}\)
and the corresponding near-degenerate batch family
\(\mathfrak B_{\mathrm{nd}}(t):=\left\{ B\subseteq \mathcal N_3(t): r_{\max}(t)\in B,\quad 1\le |B|\le B_{\max}(t) \right\}\).
Define a hard batch-structural gate
\(\Gamma_{\mathrm{batch}}^{\mathrm{struct}}(B;t)\in\{0,1\}\) to mean
that the records in \(B\) are jointly insertable on the current
branch/state, satisfy all route-declared hard cross-record constraints,
and determine a well-defined batched scaffold update
\(\mathcal O_t\oplus B\) in the fixed canonical insertion order. This is
an evaluability predicate for the route, not a theorem-level commutation
predicate: a conservative TETRIS-style route may instantiate it with
pairwise exact commutation and non-overlapping support, while the
route-neutral reduced-plane law only requires the proposed batch to be
well-defined before \(G_B(t)\), \(\widetilde H_B(t)\), and
\(\Delta E_B^{\mathrm{TR}}(t)\) are evaluated. The structurally feasible
batch family is
\(\mathfrak B_{\mathrm{struct}}(t):=\left\{ B\subseteq \mathcal S_3(t): \Gamma_{\mathrm{batch}}^{\mathrm{struct}}(B;t)=1 \right\}\).
A controller-side transient beam fallback may also be defined on the
same rerank
shell.\footnote{One convenient rule is to set $\mathcal T_{\mathrm{tie}}(t):=\{\,r\in\mathcal S_3(t): s_t(r)\ge \max(\eta_{\mathrm{tie}}\,s_t(r_{\max}(t)),\,s_t(r_{\max}(t))-\delta_{\mathrm{tie}})\,\}$. If $|\mathcal T_{\mathrm{tie}}(t)|>1$ and the controller elects not to resolve the shell by an immediate one-scaffold batch, it may retain a highest-scoring subset $\mathcal A_{\mathrm{tie}}(t)\subseteq \mathcal T_{\mathrm{tie}}(t)$ with $1\le |\mathcal A_{\mathrm{tie}}(t)|\le B_{\mathrm{tie}}(t)$ and pass one retained alternative per record into the scaffold-beam machinery of \S 11.5.2.}

Common-window reduced geometry: write the current state as
\(|\psi_t\rangle:=|\psi(\mathcal O_t,\theta_t)\rangle, \qquad Q_{\psi_t}:=I-|\psi_t\rangle\langle \psi_t|\),
where \(Q_{\psi_t}\) is the horizontal projector. For any batch
\(B\in \mathfrak B_{\mathrm{struct}}(t)\), define the common inherited
refit window \(W_B(t):=\bigcup_{r\in B} W_r(t)\) and the corresponding
inherited horizontal tangent matrix
\(T_B(t):=[\,t_j(t)\,]_{j\in W_B(t)}, \qquad t_j(t):=Q_{\psi_t}\,\partial_{\theta_j}|\psi_t\rangle\).
Let \(E_t(\alpha,\delta\theta_{W_B};B)\) denote the local energy
obtained by inserting the records in \(B\) with batch amplitudes
\(\alpha=(\alpha_r)_{r\in B}\) and varying the inherited coordinates on
\(W_B(t)\) by \(\delta\theta_{W_B}\). For each \(r\in B\), let
\(|\psi_{t,r}(\alpha_r)\rangle\) denote the state obtained by inserting
only record \(r\) at amplitude \(\alpha_r\) into the current scaffold
while keeping inherited coordinates fixed. Its zero-initialized
horizontal candidate tangent is
\(d_r(t):=Q_{\psi_t}\,\partial_{\alpha_r}|\psi_{t,r}(\alpha_r)\rangle\Big|_{\alpha_r=0}\).
Let \(\mathsf H_{W_B}(t)\) denote the inherited-window Hessian block of
\(E_t(\alpha,\delta\theta_{W_B};B)\) with respect to the coordinates on
\(W_B(t)\) at \((\alpha,\delta\theta_{W_B})=(0,0)\), and set
\(M_B(t):=\frac12\Bigl(\mathsf H_{W_B}(t)+\mathsf H_{W_B}(t)^\top\Bigr)+\lambda_H I, \qquad \lambda_H>0.\)
For each \(r\in B\), define the mixed candidate-window curvature vector
by
\(\bigl(b_r^{(B)}(t)\bigr)_j:=\left.\partial_{\alpha_r}\partial_{\theta_j} E_t(\alpha,\delta\theta_{W_B};B)\right|_{(\alpha,\delta\theta_{W_B})=(0,0)}, \qquad j\in W_B(t)\).
The reduced candidate tangent after linear relaxation of the common
inherited window is then
\(u_r^{(B)}(t):=d_r(t)-T_B(t)\,M_B(t)^{-1}b_r^{(B)}(t)\), hence the
effective batch plane is
\(\Pi_B(t):=\operatorname{span}\{u_r^{(B)}(t):r\in B\}\subseteq T_{\psi_t}\mathbb{P}(\mathcal H)\).
Its reduced Gram matrix is
\(\bigl(G_B(t)\bigr)_{rs}:=\Re\langle u_r^{(B)}(t),u_s^{(B)}(t)\rangle, \qquad r,s\in B\),
and in particular \(\bigl(G_B(t)\bigr)_{rr}=F_r^{\mathrm{red},(B)}(t)\)
is the reduced norm of record \(r\) in the common batch reduction. Next
define the symmetrized raw candidate-candidate curvature on the same
common reduction by
\(h_{rs}^{(B)}(t):=\frac12 \left. \bigl( \partial_{\alpha_r}\partial_{\alpha_s} + \partial_{\alpha_s}\partial_{\alpha_r} \bigr) E_t(\alpha,\delta\theta_{W_B};B) \right|_{(\alpha,\delta\theta_{W_B})=(0,0)}, \qquad r,s\in B,\)
and then the reduced curvature matrix
\(\bigl(\widetilde H_B(t)\bigr)_{rs}:=h_{rs}^{(B)}(t)-\bigl(b_r^{(B)}(t)\bigr)^\top M_B(t)^{-1} b_s^{(B)}(t)\).
For \(r=s\), the entry \(\bigl(\widetilde H_B(t)\bigr)_{rr}\) is the
usual single-record reduced curvature in the common batch window; for
\(r\neq s\), the entry \(\bigl(\widetilde H_B(t)\bigr)_{rs}\) is the
mixed second-order coupling after inherited-window relaxation. The
diagnostic coherences are
\(\mu_{\mathrm{tan}}(B;t):=\max_{r\neq s\in B} \frac{|(G_B(t))_{rs}|} {\sqrt{(G_B(t))_{rr}(G_B(t))_{ss}}+\varepsilon}\)
and
\(\mu_{\mathrm{curv}}(B;t):=\max_{r\neq s\in B} \frac{|(\widetilde H_B(t))_{rs}|} {\sqrt{\bigl((\widetilde H_B(t))_{rr}\bigr)^{(+)} \bigl((\widetilde H_B(t))_{ss}\bigr)^{(+)} }+\varepsilon}, \qquad x^{(+)}:=\max\{x,0\}\);
these may be used as optional early pruning diagnostics, but the actual
batch objective is the joint reduced trust-region gain below.

Joint reduced gain and additivity: let
\(g_{B,\mathrm{lcb}}(t):=\bigl(g_{\mathrm{lcb}}(r;t)\bigr)_{r\in B}\),
and let \(\bigl(\widetilde H_B(t)\bigr)^{(+)}\) denote the
positive-semidefinite part of \(\widetilde H_B(t)\). The batch
trust-region objective below is therefore a conservative convexified
local gain model on the reduced batch plane: \[
\Delta E_B^{\mathrm{TR}}(t)
:=
\max_{\alpha\in\mathbb R^{|B|},\ \alpha^\top G_B(t)\alpha\le \rho_B(t)^2}
\left[
g_{B,\mathrm{lcb}}(t)^\top \alpha
-
\frac12 \alpha^\top \bigl(\widetilde H_B(t)\bigr)^{(+)}\alpha
\right].
\] Here \(\rho_B(t)>0\) is the Phase-3 batch trust-region radius. If
record identity already includes the favorable orientation, one may
additionally impose \(\alpha\ge 0\) componentwise. For each \(r\in B\),
define the contextual single-record gain in the same common reduction by
\(\Delta E_{r\mid B}^{\mathrm{TR}}(t):=\max_{\beta\in\mathbb R,\ \beta^2 (G_B(t))_{rr}\le \rho_B(t)^2} \left[ g_{\mathrm{lcb}}(r;t)\,\beta - \frac12 \bigl((\widetilde H_B(t))_{rr}\bigr)^{(+)}\beta^2 \right].\)
Then the additivity defect is
\(\delta_{\mathrm{add}}(B;t):=\left[ 1- \frac{\Delta E_B^{\mathrm{TR}}(t)} {\sum_{r\in B}\Delta E_{r\mid B}^{\mathrm{TR}}(t)+\varepsilon} \right]_{+}.\)
Small \(\delta_{\mathrm{add}}(B;t)\) means that the joint reduced model
is approximately additive on the common batch plane.

Batch choice and greedy implementation: the geometric admissible batch
family is
\(\mathfrak B_{\mathrm{geom}}(t):=\left\{ B\in \mathfrak B_{\mathrm{nd}}(t)\cap \mathfrak B_{\mathrm{struct}}(t): \lambda_{\min}\!\bigl(G_B(t)\bigr) \ge \tau_{\mathrm{rank}}\, \frac{\operatorname{tr}G_B(t)}{|B|}, \quad \delta_{\mathrm{add}}(B;t)\le \tau_{\mathrm{add}} \right\}\).
When a more aggressive pruning rule is desired, one may additionally
require
\(\mu_{\mathrm{tan}}(B;t)\le \tau_{\mathrm{tan}}, \qquad \mu_{\mathrm{curv}}(B;t)\le \tau_{\mathrm{curv}}.\)
The principled post-shortlist batch choice is \[
B_t^\star
\in
\arg\max_{B\in \mathfrak B_{\mathrm{geom}}(t)}
\Delta E_B^{\mathrm{TR}}(t).
\] If \(\mathfrak B_{\mathrm{geom}}(t)=\varnothing\), no Phase-3 batch
is admitted from the current shortlist. A greedy implementation avoids
exhaustive search: initialize \(B_t^{(1)}:=\{r_{\max}(t)\}\). Given a
current batch \(B_t^{(k)}\), define the admissible add-on family
\(\mathcal A_k(t):=\left\{ r\in \mathcal N_3(t)\setminus B_t^{(k)}: B_t^{(k)}\cup\{r\}\in \mathfrak B_{\mathrm{nd}}(t)\cap \mathfrak B_{\mathrm{struct}}(t) \right\}\).
For each \(r\in \mathcal A_k(t)\), define the marginal joint gain
\(\Delta_{\mathrm{marg}}(r\mid B_t^{(k)};t):=\Delta E_{B_t^{(k)}\cup\{r\}}^{\mathrm{TR}}(t)-\Delta E_{B_t^{(k)}}^{\mathrm{TR}}(t).\)
Choose
\(r_k^{\mathrm{best}}(t)\in \arg\max_{r\in \mathcal A_k(t)} \Delta_{\mathrm{marg}}(r\mid B_t^{(k)};t)\)
subject to
\(\lambda_{\min}\!\bigl(G_{B_t^{(k)}\cup\{r\}}(t)\bigr) \ge \tau_{\mathrm{rank}}\, \frac{\operatorname{tr}G_{B_t^{(k)}\cup\{r\}}(t)}{|B_t^{(k)}|+1}, \qquad \delta_{\mathrm{add}}(B_t^{(k)}\cup\{r\};t)\le \tau_{\mathrm{add}},\)
and, when enabled,
\(\mu_{\mathrm{tan}}(B_t^{(k)}\cup\{r\};t)\le \tau_{\mathrm{tan}}, \qquad \mu_{\mathrm{curv}}(B_t^{(k)}\cup\{r\};t)\le \tau_{\mathrm{curv}}.\)
If no such \(r\) exists, or if the best feasible marginal gain is
nonpositive, stop and set \(B_t^\star:=B_t^{(k)}\). Otherwise append the
best add-on, \(B_t^{(k+1)}:=B_t^{(k)}\cup\{r_k^{\mathrm{best}}(t)\}\),
and continue until no further feasible positive-gain append exists or
until \(|B_t^{(k)}|=B_{\max}(t)\). The deployed Phase-3 action is the
final batch \(B_t^\star\). If \(|B_t^\star|=1\), this reduces to the
usual single-record admission. If \(|B_t^\star|>1\), the controller
inserts the batch on the current branch in the canonical batch order and
performs the common-window refit on \(W_{B_t^\star}(t)\).

\subsubsection{11.4.5 Planned size-aware Hamiltonian-agnostic Phase-3
policy
optimization}\label{planned-size-aware-hamiltonian-agnostic-phase-3-policy-optimization}

The planned policy-optimization layer treats the Phase-3 selector above
as a controlled decision system, not as a Hamiltonian-label lookup
table. Let \(h\in\mathfrak H\) denote a Hamiltonian instance, \(L\) its
size, \(\mathcal G_h(L)\) the resolved candidate-generator pool, and
\(\xi_h\) a transferable descriptor vector containing scale-normalized
couplings, graph/locality summaries, sector data, term-family counts,
operator-family features, and algebraic graph descriptors. The algebraic
portion of \(\xi_h\) includes local commutation degree histograms, local
noncommutation degree histograms, commuting-cluster size distribution,
noncommuting-bridge distribution, flat-cluster saturation statistics,
support-overlap graph statistics, and per-family commutation profiles.
For a fixed pair \((h,L)\), the instance oracle policy is \[
\phi^\star_{h,L}
\in
\arg\min_{\phi\in\Phi_L} J_{h,L}(\phi),
\qquad
J_{h,L}(\phi):=
\mathbb E_{\omega\sim\Omega_{h,L}}\!\left[\ell_{h,L,\omega}(\phi)\right]
+\lambda_{\mathrm{fail}}\,\mathbb P_{\omega}\!\left(F_{h,L,\omega}(\phi)\right),
\] where \(\omega\) ranges over seeds, shot/noise draws, and admissible
initialization choices; \(\ell\) is the terminal cost such as final
energy error plus normalized measurement and compile burden; and \(F\)
is a run failure event such as nonconvergence, invalid scaffold growth,
or violation of a hard controller contract. This oracle is an analysis
object: it is allowed to specialize to the held pair \((h,L)\), so it
gives the best per-instance policy one could hope to match.

For heavily bosonic benchmark families, the small-system training suite
should not use a single-excitation oscillator cutoff as the only
evidence surface. Let \(n_{b}^{\max}\) denote the local boson or phonon
cutoff. The default bosonic stress lane uses \(n_{b}^{\max}=2\) for the
pure boson-chain and spin-boson families, while retaining cheaper
\(n_{b}^{\max}=1\) lanes where they are intentionally used as smoke
tests or HH-specific historical baselines. This prevents the outer
policy from fitting selector weights to an artificially two-level
oscillator and exposes whether shortlisting, shot-cost penalties,
insertion, beam, and rollback still behave when each bosonic site has
three local Fock levels.

The same suite should include the closed-shell molecular route as an
electronic anchor, even though its Hamiltonian construction uses a
molecular-integral JSON rather than a lattice-coupling tuple. The
molecular lane uses explicit roles: H2/STO-3G remains the lightweight
control, LiH/STO-3G is the primary molecular train anchor, and
H2O/STO-3G remains the transfer/stress molecule. Each molecular
benchmark samples the same selector policy vector \(\phi\), so this lane
tests whether the learned shortlist, cost, beam, insertion, and rollback
weights transfer beyond lattice-generated operator families while
keeping the legacy H2/H2O benchmark identities stable.

The deployable target is instead a single Hamiltonian-agnostic,
size-aware map \[
\phi_{\mathrm{agn}}:\ (L,\xi_h)\longmapsto
\Bigl(\tau_k,\eta_k,N_k,B_{\max},\rho_B,\lambda_H,
\{B_{1,\ell},B_{2,\ell},c_\ell^{\mathrm{abs}},c_\ell^{\mathrm{rel}},
\eta_{1,\ell},\eta_{2,\ell}^{\mathrm{in}}\}_{\ell\in\mathfrak L},
\beta_{\mathrm{comm\_recover}},\beta_{\mathrm{curv\_recover}},
\beta_{\mathrm{bridge\_risk}},
\text{family lane quotas}\Bigr)_{k=1}^{3},
\] which emits Phase-1/2/3 thresholds, lane budgets, lane-survival
patience factors, batch caps, trust radii, regularizers, and
prune-recovery risk weights before the online selector is evaluated.
Algebraic descriptors in \(\xi_h\) may tune lane exposure, lane
patience, and prune-recovery caution; they do not create score bonuses
for \(S_1\), \(S_2\), or \(S_3\). Its population objective should prefer
good average behavior while protecting the worst practical cases: \[
\begin{aligned}
\mathcal J_{\mathrm{agn}}(\phi_{\mathrm{agn}})
&:=
\mathbb E_{(h,L)\sim\mathcal D}\left[
\mathbb E_{\omega}\,\ell_{h,L,\omega}\!\left(\phi_{\mathrm{agn}}(L,\xi_h)\right)
\right] \\
&\quad+
\lambda_{\mathrm{cvar}}\,
\operatorname{CVaR}_{\alpha}\!\left(
\ell_{h,L,\omega}\!\left(\phi_{\mathrm{agn}}(L,\xi_h)\right)
\right)
+\lambda_{\mathrm{rob}}\,
\sup_{(h,L)\in\mathcal U}\left[
J_{h,L}\!\left(\phi_{\mathrm{agn}}(L,\xi_h)\right)-J^{\mathrm{base}}_{h,L}
\right]_{+} \\
&\quad+
\lambda_{\mathrm{fail}}\,
\mathbb P_{(h,L),\omega}\!\left(
F_{h,L,\omega}\!\left(\phi_{\mathrm{agn}}(L,\xi_h)\right)
\right).
\end{aligned}
\] The first term optimizes mean performance; the CVaR term penalizes
the high-loss tail; the robust term prevents a small adversarial or
withheld family \(\mathcal U\) from being sacrificed; and the failure
term keeps apparently cheap policies from winning by producing brittle
or invalid runs.

The natural measure of success is the oracle gap \[
\Delta_{\mathrm{oracle}}(h,L;\phi_{\mathrm{agn}})
:=
J_{h,L}\!\left(\phi_{\mathrm{agn}}(L,\xi_h)\right)
-
J_{h,L}(\phi^\star_{h,L})
\ge 0,
\] reported both as a mean over \(\mathcal D\) and as a tail statistic
over withheld Hamiltonians and larger sizes. A small gap means the
agnostic policy has learned size and family regularities rather than
memorizing a particular Hamiltonian.

Shortlist and pool budgets should scale explicitly with both size and
available generator count. A compact family of caps is \[
N_k(L,|\mathcal G_h(L)|;\zeta_k)
:=
\min\!\left\{
|\mathcal G_h(L)|,
\left\lceil a_k L^{p_k}\bigl[\log(1+|\mathcal G_h(L)|)\bigr]^{q_k}\right\rceil
\right\},
\qquad k\in\{1,2,3\},
\] with learned or scheduled \(\zeta_k=(a_k,p_k,q_k)\), constrained so
that \(N_3\le N_2\le N_1\) after rounding. Thus larger systems may
receive larger absolute candidate budgets, while the fraction of the
full pool inspected can still shrink when \(|\mathcal G_h(L)|\) grows
faster than the useful front.

These shortlists are intended to be active control surfaces, not
disabled compatibility knobs. A concrete tuning layer should therefore
tune \(N_1,N_2\), the Phase-2 shortlist fraction, the insertion-position
probe cap \(M_{\mathrm{probe}}\), and the cost weights that trade energy
proxy against compile and measurement burden. The goal is to reduce shot
and compile cost without collapsing expressivity; broad-pool
expressivity is preserved by scaling the budgets with
\((L,|\mathcal G_h(L)|)\) and by retaining beam, insertion, rollback,
batch, and prune alternatives at the controller level.

Operator-family controls enter only as transferable lane-budget priors
and quotas. If \(f(r)\) is the abstract family of record \(r\), the
policy may emit smooth family-lane quotas \(q_{\ell,f}(L,\xi_h,t)\) and
family-lane exposure weights \(\pi_{\ell,f}(L,\xi_h,t)\). These
quantities can change how many Phase-1 or Phase-2 seats a family
receives inside a lane, for example by constraining \[
\left|\{r\in\mathcal S_{1,\ell}(t):f(r)=f\}\right|
\le q_{\ell,f}(L,\xi_h,t),
\] or by adjusting \(B_{1,\ell}(t)\) and \(B_{2,\ell}(t)\) before
records are ranked. They may say that displacement-like, hopping-like,
squeeze-like, or Hubbard-like families tend to deserve more or less
exposure in a size/coupling regime. They do not alter the geometric
score of any record and must not become exact label masks such as
\(\mathbf 1\{m\in\mathcal M^{\mathrm{best}}_{h,L}\}\), because such
masks would encode the answer for one Hamiltonian instead of a
transferable policy.

The final validation should freeze \(\phi_{\mathrm{agn}}\) after
training or tuning on smaller and mixed-family surfaces, then evaluate
it on held-out \(L=4\) and larger Hamiltonians. The report should
compare mean loss, CVaR loss, failure rate, budget use
\(N_k(L,|\mathcal G|)\), and oracle gap against the per-\((h,L)\) oracle
\(\phi^\star_{h,L}\), a size-blind baseline, and any label-specific
baseline. Passing this validation means the policy scales through
\(L=4+\) by using size-aware and family-transferable structure, not by
recognizing exact operator labels from the training set.

\subsection{\texorpdfstring{11.5 Post \(m_p\) Application Entails
Windowed-Refitting}{11.5 Post m\_p Application Entails Windowed-Refitting}}\label{post-m_p-application-entails-windowed-refitting}

In ADAPT-VQE the variational problem is
\(\min_{\boldsymbol\theta} E_L(\boldsymbol\theta)\) with
\(E_L(\boldsymbol\theta)=\langle\phi_0|(\prod_{\ell=L}^{1}e^{i\theta_\ell G_\ell})H(\prod_{\ell=1}^{L}e^{-i\theta_\ell G_\ell})|\phi_0\rangle\),
and the measurement bottleneck arises because one iteration must
estimate not only \(E_L\) but also pool-selection scores
\(s_\mu(\boldsymbol\theta)=|\langle\phi_0|(\prod_{\ell=L}^{1}e^{i\theta_\ell G_\ell}),i[H,A_\mu],(\prod_{\ell=1}^{L}e^{-i\theta_\ell G_\ell})|\phi_0\rangle|\)
for \(\mu\in\mathcal P\), reoptimization gradients
\(\partial_{\theta_j}E_L(\boldsymbol\theta)\), and, for geometry-aware
updates, metric entries
\(M_{jk}(\boldsymbol\theta)=\Re\langle \partial_{\theta_j}[(\prod_{\ell=1}^{L}e^{-i\theta_\ell G_\ell})|\phi_0\rangle]|\partial_{\theta_k}[(\prod_{\ell=1}^{L}e^{-i\theta_\ell G_\ell})|\phi_0\rangle]\rangle\),
so at precisions \(\epsilon_s,\epsilon_g,\epsilon_M\) the shot cost
obeys
\(N_{\mathrm{iter}}\gtrsim \sum_{\mu\in\mathcal P}\mathrm{Var}(\hat s_\mu)/\epsilon_s^2+\sum_{j=1}^{L}\mathrm{Var}(\widehat{\partial_{\theta_j}E_L})/\epsilon_g^2+\sum_{j,k=1}^{L}\mathrm{Var}(\hat M_{jk})/\epsilon_M^2\);
the proposed strategy is to choose disjoint sets
\(F,A\subset{1,\dots,L}\) with \(F\cup A={1,\dots,L}\), impose
\(\dot\theta_j=0\) for \(j\in F\), and optimize only \(j\in A\), which
reduces the reparameterization sector to
\(\sum_{j\in A}\mathrm{Var}(\widehat{\partial_{\theta_j}E_L})/\epsilon_g^2+\sum_{j,k\in A}\mathrm{Var}(\hat M_{jk})/\epsilon_M^2\)
and therefore removes the growth coming from repeatedly updating old
coordinates, but it leaves the pool-screening term
\(\sum_{\mu\in\mathcal P}\mathrm{Var}(\hat s_\mu)/\epsilon_s^2\)
unchanged and it does not shorten the physical circuit, since every shot
still executes the full depth-\(L\) scaffold with entangling-gate count
\(N_{2q}(L)\) and runtime \(t_{\mathrm{circ}}(L)\), so hardware fidelity
still degrades roughly as
\(e^{-t_{\mathrm{circ}}(L)/T_1}e^{-t_{\mathrm{circ}}(L)/T_2}(1-\varepsilon_{2q})^{N_{2q}(L)}\);
consequently the freezing strategy cures only the reoptimization
component of the ADAPT bottleneck, and the limiting cost after that is
either the unchanged operator-pool screening overhead or, if screening
is separately controlled, the conventional VQE bottleneck of resolving
\(E_L\) on a deep noisy ansatz, with shots scaling as
\(N_E\gtrsim \mathrm{Var}(H_{\mathrm{meas}}\mid \prod_{\ell=1}^{L}e^{-i\theta_\ell G_\ell},\phi_0,\text{noise})/(\Delta E)^2\)
to distinguish an energy change \(\Delta E\).

Let the length-\((m)\) scaffold be
\(U^{(m)}(\boldsymbol{\theta})=\prod_{j=1}^{m} e^{-i\theta_j G_j}\),
where \((G_j)\) is the \((j)\)-th appended generator and \((\theta_j)\)
its coefficient. The reparameterization motif is to optimize only an
active subset \((\mathcal A_m\subseteq {1,\dots,m})\), chosen by an
activity functional
\(a_j^{(m)}:=\lambda_{\mathrm{rec}}\,r_j^{(m)}+\lambda_{\mathrm{mob}}\,v_j^{(m)}-\lambda_{\mathrm{age}}\,\phi(m-j)-\lambda_{\mathrm{stag}}\,s_j^{(m)}\),
and then set \((\mathcal A_m=\{j: a_j^{(m)}>\tau_m\})\). Here
\((r_j^{(m)})\) is a recency weight, large for recently appended
generators;
\((v_j^{(m)}:=\sum_{k=j}^{m-1}\lvert \Delta\theta_j^{(k)}\rvert)\) is a
lifetime mobility measure, large for parameters that have historically
moved substantially; \((\phi(m-j))\) is an age penalty, biasing against
early-added generators; and \((s_j^{(m)})\) is a stagnation indicator,
large when a parameter has remained nearly inert over recent refits.
Thus the inclusion bias is toward recent and historically labile
generators, while the exclusion bias is toward old and typically
stagnant generators.

This is the same scaffold language used above:
\(U^{(m)}(\boldsymbol\theta)\) is the scaffold \(\mathcal O\), the
active set \(\mathcal A_m\) plays the role of the inherited refit window
\(W(r;t)\), the frozen set is \(\{1,\ldots,n+1\}\setminus W(r;t)\), and
the activity functional \(a_j^{(m)}\) is the same maturity-facing
bookkeeping later summarized by the eligible set \(\mathcal M_t\).

Concomitantly, as scaffold maturity increases, the controller shot laws
\(N_{\mathrm{shot},k}^{\mathrm{eff}}(t)\) should preserve decision-level
signal-to-noise ratio on the active block while respecting the live
phase mask. Formally, if \(\delta_{t,k}\) denotes a typical useful
signal scale among the phase-\(k\) unfrozen directions and
\(\sigma_{t,k}\) the corresponding estimator spread, then the capped
schedule-plus-SNR law above targets \[
\mathrm{SNR}_{t,k}\sim
\frac{\delta_{t,k}\sqrt{N_{\mathrm{shot},k}^{\mathrm{eff}}(t)}}{\sigma_{t,k}}
\ge \kappa_k
\] for every live phase \(k\), while a nulled phase has
\(N_{\mathrm{shot},k}^{\mathrm{eff}}(t)=0\). Hence the control law is:
as scaffold maturity grows, shrink the reoptimized subspace by freezing
archaic or stagnant directions, increase shots on the surviving live
active directions, and eventually null expensive geometric phases when
the runway forecast says the scaffold is mature enough.

\subsubsection{11.5.1 Branch-local ADAPT prune-recoverability ladder:
nomination, refit, and
rollback}\label{branch-local-adapt-prune-recoverability-ladder-nomination-refit-and-rollback}

All quantities in this subsection are branch-local. On a single active
branch we suppress the branch label. Let
\((\mathcal O_t^-,\theta_t^-,E_t^-)\) be the committed branch state
immediately before the accepted Phase-3 payload \(\mathcal A_t^\star\)
is applied, and let \((\mathcal O_t^+,\theta_t^+,E_t^+)\) be the
committed post-admission state after the ordinary local refit on that
same branch has completed. The admitted-step gain is \[
\Delta_{\mathrm{adm}}(t):=E_t^- - E_t^+,
\qquad
\mathcal J_t:=\{1,\dots,n_t^+\}.
\] Each committed coordinate \(j\in\mathcal J_t\) carries the
operational metadata \[
(m_j,\iota_j,r_j,\mathfrak c_j,\kappa_j,c_j(t),\{a_j(q)\}_{q\le t},
\operatorname{noncomm\_bridge\_score}_j,
\operatorname{recovery\_class}_j),
\] where \(m_j\) is the generator label, \(\iota_j\) is the branch-local
admission step, \(r_j\) is the admission record, \(\mathfrak c_j\) is
the generator-family or motif class, \(\kappa_j:=K_3(r_j;\iota_j)\) is
the stored selector burden, \(c_j(t)\) is the prune-cooldown counter,
\(a_j(q)\) is the transported normalized amplitude history,
\(\operatorname{noncomm\_bridge\_score}_j\) is the static bridge-risk
telemetry, and \(\operatorname{recovery\_class}_j\) is the auditable
prune outcome. Optional algebraic telemetry such as commuting-cluster
labels, local commutation degrees, and flat-cluster saturation may be
stored with the coordinate, but it is not part of the main deletion
authority below.

The normalized amplitude history is \[
a_j(q):=
\frac{\left|\operatorname{wrap}_{\mathfrak c_j}\theta_j(q)\right|}
{s_{\mathfrak c_j}+\varepsilon_s},
\] using the principal-angle wrap for periodic rotation coordinates and
the identity for nonperiodic coordinates. The recovery class is
initially unset, \[
\operatorname{recovery\_class}_j\in
\{\varnothing,\mathrm{flat\_redundant},\mathrm{curvature\_compensated},\mathrm{failed}\},
\] and is set only after a prune trial.

The target object is not whether coordinate \(j\) is small; it is
whether the scaffold with \(j\) deleted can recover the same variational
energy after the surviving coordinates compensate. With
\(\mathcal O_{t,-j}^+\) denoting the scaffold obtained by deleting
\(j\), define \[
\boxed{
\delta E_j^\star(t)
:=
\inf_{\xi\in\Theta_{\mathcal O_{t,-j}^+}}
E(\mathcal O_{t,-j}^+,\xi)
-
E(\mathcal O_t^+,\theta_t^+).
}
\] Small amplitude is only one sufficient-looking nomination signal, \[
|\theta_{t,j}|\approx0
\Longrightarrow
\delta E_j^\star(t)\approx0,
\qquad
\delta E_j^\star(t)\approx0
\not\Longrightarrow
|\theta_{t,j}|\approx0.
\] Thus prune pressure may increase opportunity, candidate breadth, and
remove-refit breadth, but it must not weaken final safety. The
branch-local prune chain is \[
\boxed{
\mathcal J_t
\to
\mathcal M_t
\to
\mathcal N_1(t)
\to
\{W_j^{[s]}(t)\}_{s=0}^{4}
\to
\Gamma_j^{\mathrm{proj},[s]}(t)
\to
\mathcal C_{\mathrm{trial}}(t)
\to
\mathsf D_j^{[s,N]}
\to
\operatorname{commit}/\operatorname{rollback}.
}
\] Here \(\mathcal M_t\) is the mature deletion-target universe,
\(\mathcal N_1(t)\) is the cheap deletion-target nomination set,
\(W_j^{[s]}(t)\) are surviving-coordinate refit windows for a fixed
target \(j\), \(\Gamma_j^{\mathrm{proj},[s]}\) is an optional cached
geometry screen, \(\mathcal C_{\mathrm{trial}}(t)\) is the capped set
sent to exact remove-refit, and \(\mathsf D_j^{[s,N]}\) is the
confidence-controlled decision. No amplitude score, commutation
relation, cached projection, or Schur surrogate deletes a coordinate by
itself.

\paragraph{Prune 0: permission, maturity, and protected
coordinates}\label{prune-0-permission-maturity-and-protected-coordinates}

Prune may open after an accepted admission, after a no-admit plateau, at
scheduled compression checkpoints, or at the terminal scaffold. Let
\(\mathcal T_{\mathrm{chk}}\subseteq\mathbb N\) be the scheduled
checkpoint set, let \(\operatorname{Plateau}_t\) mean that no useful
admission has survived the current controller patience window, and let
\(\operatorname{StableRefit}_t\) mean that the latest local refit
completed without optimizer or numerical failure. The permission gate is
\[
\begin{aligned}
\Gamma_{\mathrm{prune}}^{\mathrm{perm}}(t)
&:=
\mathbf 1[|\mathcal O_t^+|\ge 2]\,
\mathbf 1[\operatorname{StableRefit}_t]\,
\mathbf 1[u_{\mathrm{sat}}(t)\ge\tau_{\mathrm{med}}]\\
&\quad\times
\mathbf 1[
\text{accepted admission at }t
\ \vee\
\operatorname{Plateau}_t
\ \vee\
t\in\mathcal T_{\mathrm{chk}}
\ \vee\
t=t_{\mathrm{terminal}}
].
\end{aligned}
\] The scale used later for prune tolerance is \[
\Delta_{\mathrm{scr}}(t)
:=
\max\!\bigl(
\tau_{\mathrm{use},t},
\operatorname{EWMA}_{\mathrm{recent}}(\Delta_{\mathrm{adm}}),
\Delta_{\mathrm{num}}
\bigr).
\] A scalar prune-pressure signal may schedule candidate caps, window
fractions, and cached-screen thresholds. For example, with normalized
pressure \(\bar\lambda_{\mathrm{pr}}(t)\in[0,1]\), one may set \[
Q_{\mathrm{pr}}(t)
:=
\left\lceil
Q_{\mathrm{pr}}^{\min}+
(Q_{\mathrm{pr}}^{\max}-Q_{\mathrm{pr}}^{\min})\bar\lambda_{\mathrm{pr}}(t)
\right\rceil,
\] \[
\rho_s(t)
:=
\operatorname{clip}_{[0,1]}\!\left(
\rho_s^{\min}+(\rho_s^{\max}-\rho_s^{\min})\bar\lambda_{\mathrm{pr}}(t)
\right),
\qquad s=1,\dots,4,
\] and may tighten or relax \(\tau_F^{\mathrm{pr}}(t)\) and
\(\tau_{\mathrm{Schur}}^{\mathrm{pr}}(t)\) according to the scheduled
compression policy. This pressure changes what is worth testing; it does
not change the final deletion authority.

Newly admitted coordinates are protected for \(\ell_{\mathrm{protect}}\)
admitted steps, \[
\mathcal P_{\mathrm{protect}}(t)
:=
\{j\in\mathcal J_t:t-\iota_j<\ell_{\mathrm{protect}}\}.
\] Let \(\operatorname{age}_{j,t}:=t-\iota_j\), and let
\(s_j^{\mathrm{stag}}(t)\) be the branch-local coordinate stagnation
indicator, with larger values meaning less recent parameter motion. The
mature deletion-target universe is \[
\boxed{
\mathcal M_t
:=
\Bigl\{
 j\in\mathcal J_t:
\operatorname{age}_{j,t}\ge a_{\min},
\ j\notin\mathcal P_{\mathrm{protect}}(t),
\ c_j(t)=0
\Bigr\}.
}
\] Prune eligibility is therefore branch-local and conservative: old
enough, not recently admitted, and not under temporary cooldown.
Staleness affects cheap nomination through a normalized feature, not
through hard eligibility.

\paragraph{Prune 1: cheap deletion-target nomination and pairwise
compensator
evidence}\label{prune-1-cheap-deletion-target-nomination-and-pairwise-compensator-evidence}

Prune 1 has two outputs. First, it nominates deletion targets that are
cheap enough to test. Second, it records pairwise evidence that Prune 2
uses to build compensator windows. Let \(\mathcal H_t\) be the recent
time-history region of accepted or refit states used for empirical
co-motion estimates. This is a time-index region, not a coordinate refit
window.

Define activity, motion, and stationarity summaries \[
\widetilde A_j(t)
:=
\frac{A_j^{\mathrm{late}}(t)}{A_{\mathrm{ref}}+\varepsilon_A},
\qquad
\widetilde V_j(t)
:=
\frac{V_j^{\mathrm{late}}(t)}{V_{\mathrm{ref}}+\varepsilon_V},
\qquad
\widetilde G_j(t)
:=
\frac{|g_j^{\mathrm{late}}(t)|}{G_{\mathrm{ref}}+\varepsilon_G}.
\] Here \(A_j^{\mathrm{late}}\) is an EWMA or windowed mean of the
normalized amplitude history, \[
A_j^{\mathrm{late}}(t)
:=
\operatorname{EWMA}_{q\in\mathcal H_t} a_j(q),
\] \(V_j^{\mathrm{late}}\) is a recent parameter-motion scale, and
\(g_j^{\mathrm{late}}\) is a cached gradient or stationarity signal.
Define normalized burden, age, and staleness features by \[
\widetilde\kappa_j(t):=\frac{\kappa_j}{\kappa_{\mathrm{ref}}+\varepsilon_\kappa},
\qquad
z_j^{\mathrm{age}}(t):=\operatorname{norm}(\operatorname{age}_{j,t}),
\qquad
z_j^{\mathrm{stale}}(t):=\operatorname{norm}(s_j^{\mathrm{stag}}(t)).
\] For \(i\ne j\), define the pairwise historical co-motion score \[
\boxed{
C_{ij}^{\mathrm{hist}}(t)
:=
\max\!\left\{
|\operatorname{corr}_{\mathcal H_t}(\theta_j,\theta_i)|,
|\operatorname{corr}_{\mathcal H_t}(g_j,g_i)|,
|\operatorname{corr}_{\mathcal H_t}(\Delta\theta_j,\Delta\theta_i)|
\right\}.
}
\] If tangent data are current and cached, define the pairwise
tangent-overlap score \[
\boxed{
C_{ij}^{\mathrm{tan}}(t)
:=
\frac{
\left|\operatorname{Re}\langle\bar\tau_j,\bar\tau_i\rangle\right|
}{
\|\bar\tau_j\|_2\,\|\bar\tau_i\|_2+\varepsilon_T
},
}
\] where \(\bar\tau_i:=Q_{\psi_t}\partial_{\theta_i}|\psi_t\rangle\) is
the horizontal tangent at the current post-admission state. If tangent
data are unavailable or stale, set \(C_{ij}^{\mathrm{tan}}(t)=0\). The
combined pairwise compensability evidence is \[
\boxed{
E_{ij}^{\mathrm{corr}}(t)
:=
\max\!\left(C_{ij}^{\mathrm{hist}}(t),C_{ij}^{\mathrm{tan}}(t)\right).
}
\] The static algebraic metadata are the support-overlap bit
\(\Omega_{m_i m_j}\), the bare-commutation bit \(\mathsf C_{m_i m_j}\),
and the exactness tag \(\mathsf q_{m_i m_j}\). Define the pairwise
algebraic evidence bits \[
\boxed{
E_{ij}^{\mathrm{flat}}
:=
\mathbf 1[\Omega_{m_i m_j}=1]\,
\mathbf 1[\mathsf C_{m_i m_j}=1]\,
\mathbf 1[\mathsf q_{m_i m_j}=\mathrm{exact}],
}
\] \[
\boxed{
E_{ij}^{\mathrm{curv}}
:=
\mathbf 1[\Omega_{m_i m_j}=1]\,
\mathbf 1[\mathsf C_{m_i m_j}=0]\,
\mathbf 1[\mathsf q_{m_i m_j}=\mathrm{exact}],
}
\] and \[
\boxed{
E_{ij}^{\mathrm{disj}}
:=
\mathbf 1[\Omega_{m_i m_j}=0]\,
\mathbf 1[\mathsf C_{m_i m_j}=1]\,
\mathbf 1[\mathsf q_{m_i m_j}=\mathrm{exact}].
}
\] The disjoint commuting bit is telemetry or batching context only. It
is not flat-sector redundancy evidence.

Aggregate features for cheap deletion-target ranking are \[
C_j^{\mathrm{comp}}(t)
:=
\max_{i\in\mathcal J_t\setminus\{j\}}E_{ij}^{\mathrm{corr}}(t),
\] \[
C_j^{\mathrm{tan}}(t)
:=
\max_{i\in\mathcal J_t\setminus\{j\}}C_{ij}^{\mathrm{tan}}(t),
\qquad
C_j^{\mathrm{comm}}(t)
:=
\max_{i\in\mathcal J_t\setminus\{j\}}E_{ij}^{\mathrm{flat}}E_{ij}^{\mathrm{corr}}(t).
\] Let
\(\mathcal J_{\mathrm{down}}(j;t)\subseteq\mathcal J_t\setminus\{j\}\)
be the recent downstream or later-in-order scaffold coordinates used to
estimate whether deleting \(j\) would remove a noncommuting bridge.
Define \[
\boxed{
B_j^{\mathrm{bridge}}(t)
:=
\operatorname{clip}_{[0,1]}\!\left(
\operatorname{noncomm\_bridge\_score}_j
+
\frac{
\sum_{i\in\mathcal J_{\mathrm{down}}(j;t)}E_{ij}^{\mathrm{curv}}
}{
|\mathcal J_{\mathrm{down}}(j;t)|+\varepsilon
}
\right).
}
\] Here \(C_j^{\mathrm{comp}}\), \(C_j^{\mathrm{tan}}\), and
\(C_j^{\mathrm{comm}}\) make deletion easier to nominate, while
\(B_j^{\mathrm{bridge}}\) protects bridge-like targets by raising their
cheap deletion score. The bridge feature is a risk feature for the
target \(j\), not a compensator-window definition.

The cheap recoverability-prior score is \[
\boxed{
\begin{aligned}
R_j^{\mathrm{cheap}}(t)
:=
&\ w_A\widetilde A_j
+w_V\widetilde V_j
+w_G\widetilde G_j
-w_C C_j^{\mathrm{comp}}
-w_T C_j^{\mathrm{tan}}
-w_{\mathrm{comm}}C_j^{\mathrm{comm}}\\
&+w_{\mathrm{bridge}}B_j^{\mathrm{bridge}}
-w_{\kappa}\widetilde\kappa_j
-w_{\mathrm{age}}z_j^{\mathrm{age}}
-w_{\mathrm{stale}}z_j^{\mathrm{stale}}.
\end{aligned}
}
\] Low \(R_j^{\mathrm{cheap}}\) means
\texttt{try\ testing\ this\ target,\textquotesingle{}\textquotesingle{}\ not}delete
this target.'\,' Define the cheap-score order by \[
j_1\preceq_{\mathrm{cheap}} j_2
\Longleftrightarrow
R_{j_1}^{\mathrm{cheap}}(t)\le R_{j_2}^{\mathrm{cheap}}(t),
\] with deterministic ties resolved by the fixed scaffold order.

The pure amplitude/activity-collapse lane is \[
\boxed{
\mathcal P_{\mathrm{amp}}(t)
:=
\left\{
 j\in\mathcal M_t:
\widetilde A_j(t)\le\tau_A^{\mathrm{pr}}(t),
\ \widetilde V_j(t)\le\tau_V^{\mathrm{pr}}(t),
\ \widetilde G_j(t)\le\tau_G^{\mathrm{pr}}(t)
\right\}.
}
\] The broader cheap recoverability-prior lane is \[
\boxed{
\mathcal P_{\mathrm{cheap}}(t)
:=
\operatorname{Bottom}_{Q_{\mathrm{pr}}(t)}
\left(
R_j^{\mathrm{cheap}}(t);
\ j\in\mathcal M_t
\right).
}
\] Prune 1 ends with the cheap deletion-target nomination set \[
\boxed{
\mathcal N_1(t)
:=
\mathcal P_{\mathrm{amp}}(t)
\cup
\mathcal P_{\mathrm{cheap}}(t).
}
\] The same pairwise evidence used to rank cheap nominees now constructs
the surviving-coordinate compensator windows.

\paragraph{Prune 2: typed compensator windows and cached Fubini--Schur
screens}\label{prune-2-typed-compensator-windows-and-cached-fubinischur-screens}

For each mature coordinate \(j\in\mathcal M_t\), define the
surviving-coordinate universe \[
I_{-j}:=\mathcal J_t\setminus\{j\},
\] identified with the coordinates of \(\mathcal O_{t,-j}^+\) after the
obvious index shift. A compensator window is a subset of \(I_{-j}\); it
contains surviving amplitudes allowed to refit after \(j\) is removed.
These windows are not nomination lanes.

The typed compensator pools are constructed from the Prune-1 pairwise
evidence: \[
\boxed{
W_{\mathrm{comm}}(j;t)
:=
\{i\in I_{-j}:E_{ij}^{\mathrm{flat}}=1\},
}
\] \[
\boxed{
W_{\mathrm{corr}}(j;t)
:=
\{i\in I_{-j}:E_{ij}^{\mathrm{corr}}(t)\ge\tau_{\mathrm{corr}}^{\mathrm{pr}}(t)\},
}
\] \[
\boxed{
W_{\mathrm{nc}}(j;t)
:=
\{i\in I_{-j}:E_{ij}^{\mathrm{curv}}=1\},
}
\] and \[
\boxed{
W_{\mathrm{age}}(j;t)
:=
\{i\in I_{-j}:\operatorname{age}_{i,t}\ge a_{\mathrm{comp}},\ i\notin\mathcal P_{\mathrm{protect}}(t)\}.
}
\] Thus \(W_{\mathrm{comm}}\) is the exact overlapping commuting pool,
\(W_{\mathrm{corr}}\) is the empirical or tangent-correlated pool,
\(W_{\mathrm{nc}}\) is the overlapping noncommuting pool, and
\(W_{\mathrm{age}}\) is the broad old-coordinate pool. Support-disjoint
commuting coordinates do not enter \(W_{\mathrm{comm}}\).

For a pool \(P(j;t)\subseteq I_{-j}\), write \(T_\rho(P,p)\) for the top
\(\lceil \rho |P|\rceil\) coordinates under priority score
\(p_{ij}(t)\), with deterministic ties. A conservative choice is \[
p_{ij}^{\mathrm{comm}}(t):=E_{ij}^{\mathrm{flat}}\bigl(\varepsilon_{\mathrm{comm}}+E_{ij}^{\mathrm{corr}}(t)\bigr),
\qquad
p_{ij}^{\mathrm{corr}}(t):=E_{ij}^{\mathrm{corr}}(t),
\] \[
p_{ij}^{\mathrm{nc}}(t):=E_{ij}^{\mathrm{curv}}E_{ij}^{\mathrm{corr}}(t),
\qquad
p_i^{\mathrm{age}}(t):=z_i^{\mathrm{age}}(t)+z_i^{\mathrm{stale}}(t).
\] Define truncated pools, for example, \[
W_{\mathrm{comm}}^{\rho}(j;t):=T_\rho(W_{\mathrm{comm}},p^{\mathrm{comm}}),
\] and analogously for \(W_{\mathrm{corr}}^\rho\),
\(W_{\mathrm{nc}}^\rho\), and \(W_{\mathrm{age}}^\rho\).

The nested refit ladder is recursive: \[
\boxed{W_j^{[0]}(t):=\varnothing,}
\] \[
\boxed{W_j^{[1]}(t):=W_{\mathrm{comm}}^{\rho_1}(j;t),}
\] \[
\boxed{W_j^{[2]}(t):=W_j^{[1]}(t)\cup W_{\mathrm{comm}}^{\rho_2}(j;t)\cup W_{\mathrm{corr}}^{\rho_2}(j;t),}
\] \[
\boxed{W_j^{[3]}(t):=W_j^{[2]}(t)\cup W_{\mathrm{nc}}^{\rho_3}(j;t),}
\] \[
\boxed{W_j^{[4]}(t):=W_j^{[3]}(t)\cup W_{\mathrm{age}}^{\rho_4}(j;t)\cup W_{\mathrm{term}}(j;t).}
\] Here \(0\le\rho_1\le\rho_2\le\rho_3\le\rho_4\le1\), and
\(W_{\mathrm{term}}(j;t)\subseteq I_{-j}\) is the optional broad
terminal or scheduled-compression window. The recursion guarantees \[
\varnothing=W_j^{[0]}(t)
\subseteq
W_j^{[1]}(t)
\subseteq
W_j^{[2]}(t)
\subseteq
W_j^{[3]}(t)
\subseteq
W_j^{[4]}(t)
\subseteq I_{-j}.
\] Rungs \(s\le2\) are frozen, flat-sector, or correlated recovery
attempts. Rungs \(s\ge3\) have opened noncommuting or broad compensators
and therefore are curvature-compensated unless the first accepted rung
occurred earlier.

Before launching an expensive exact remove-refit trial, a branch may
apply a cached Fubini--Schur screen. Let \[
\operatorname{valid}_{j}^{\mathrm{cache},[s]}(t)
:=
\operatorname{valid}_{j}^{\mathrm{cache}}(W_j^{[s]}(t);t)
\] mean that the cached tangent-overlap and Hessian/Schur blocks for
coordinate \(j\) were evaluated at the current branch state
\((\mathcal O_t^+,\theta_t^+)\) and for the same rung window
\(W_j^{[s]}(t)\). If this validity flag is zero, the projected screen
for that rung abstains; stale geometry cannot authorize a prune trial.

For nonempty \(W\subseteq I_{-j}\), define \[
\bigl(Q_j(W;t)\bigr)_{ik}
:=
\operatorname{Re}\langle \bar\tau_i,\bar\tau_k\rangle,
\qquad
\bigl(q_j(W;t)\bigr)_i
:=
\operatorname{Re}\langle \bar\tau_i,\bar\tau_j\rangle,
\] for \(i,k\in W\), and \[
F_j(t):=\operatorname{Re}\langle\bar\tau_j,\bar\tau_j\rangle+\varepsilon_F.
\] With \(\lambda_F>0\), the scale-free tangent residual is \[
\boxed{
\rho_j^{\mathrm{tan}}(W;t)
:=
\operatorname{clip}_{[0,1]}\!\left(
1-
\frac{
q_j(W;t)^\top\bigl(Q_j(W;t)+\lambda_F I\bigr)^{-1}q_j(W;t)
}{F_j(t)}
\right).
}
\] The amplitude-scaled Fubini deletion residual is obtained from \[
\delta_{W,F}^{\star}(j;t)
:=
\theta_{t,j}\bigl(Q_j(W;t)+\lambda_F I\bigr)^{-1}q_j(W;t),
\] \[
\boxed{
D_{j\mid W,F}^{2}(t)
:=
\theta_{t,j}^2F_j(t)
-2\theta_{t,j}q_j(W;t)^\top\delta_{W,F}^{\star}(j;t)
+
\bigl(\delta_{W,F}^{\star}(j;t)\bigr)^\top Q_j(W;t)\delta_{W,F}^{\star}(j;t).
}
\] In the unregularized ideal case, \[
D_{j\mid W,F}^{2}(t)
=
\theta_{t,j}^2
\left(F_j(t)-q_j(W;t)^\top Q_j(W;t)^+q_j(W;t)\right).
\] Thus \(\rho_j^{\mathrm{tan}}\) is a scale-free tangent-rank
diagnostic, while \(D_{j\mid W,F}^{2}\) is the state-space prune-screen
quantity.

For the energy-side screen, let \(H_W(t)\) be the Hessian block on
\(W\), let \(h_{Wj}(t)\) be the mixed Hessian vector between \(W\) and
coordinate \(j\), let \(h_{jW}(t):=h_{Wj}(t)^\top\), let
\(g_W(t):=(g_i(t))_{i\in W}\), and let \(h_{jj}(t)\) be the
one-coordinate curvature of \(j\). Define \[
R_j(W;t):=\frac12\left(H_W(t)+H_W(t)^\top\right)+\lambda_H I.
\] The gradient-aware quadratic deletion-refit surrogate is \[
\boxed{
\widehat L_j^{\mathrm{Schur}}(W;t)
:=
-g_j(t)\theta_{t,j}
+\frac12h_{jj}(t)\theta_{t,j}^2
-\frac12
\bigl(g_W(t)-h_{Wj}(t)\theta_{t,j}\bigr)^\top
R_j(W;t)^{-1}
\bigl(g_W(t)-h_{Wj}(t)\theta_{t,j}\bigr).
}
\] Near a locally refit state, the curvature-only version is \[
\widetilde h_j^{(-)}(W;t)
:=
h_{jj}(t)-h_{jW}(t)R_j(W;t)^{-1}h_{Wj}(t),
\] \[
\boxed{
L_j^{\mathrm{Schur}}(W;t)
:=
\frac12\theta_{t,j}^2\max\{\widetilde h_j^{(-)}(W;t),0\}.
}
\] If the branch is not near a stable local refit,
\(\widehat L_j^{\mathrm{Schur}}\) is the safer nomination proxy; if the
gradient/Hessian cache is missing or stale, the cached screen abstains
rather than launching new measurements only for nomination.

For \(s\in\{1,2,3,4\}\), define the rungwise projected screen \[
\boxed{
\Gamma_j^{\mathrm{proj},[s]}(t)
:=
\mathbf 1[\operatorname{valid}_{j}^{\mathrm{cache},[s]}(t)=1]\,
\mathbf 1[D_{j\mid W_j^{[s]},F}^{2}(t)\le(\tau_F^{\mathrm{pr}}(t))^2]\,
\mathbf 1[L_j^{\mathrm{Schur}}(W_j^{[s]};t)\le\tau_{\mathrm{Schur}}^{\mathrm{pr}}(t)].
}
\] Set \(\Gamma_j^{\mathrm{proj},[s]}(t)=0\) when
\(W_j^{[s]}(t)=\varnothing\). The first geometrically passing rung is \[
\boxed{
s_{\mathrm{geo}}(j;t)
:=
\min\{s\in\{1,2,3,4\}:\Gamma_j^{\mathrm{proj},[s]}(t)=1\},
}
\] with \(s_{\mathrm{geo}}(j;t)=\infty\) if no cached rung passes. The
projected flat and curvature-compensated sets are \[
\boxed{
\mathcal C_{\mathrm{flat}}^{\mathrm{proj}}(t)
:=
\{j\in\mathcal M_t:s_{\mathrm{geo}}(j;t)\le2\},
}
\] \[
\boxed{
\mathcal C_{\mathrm{curv}}^{\mathrm{proj}}(t)
:=
\{j\in\mathcal M_t:3\le s_{\mathrm{geo}}(j;t)\le4\}.
}
\] If only one side of the cache is available, weak geometry lanes may
also be recorded: \[
\mathcal P_F(t)
:=
\{j\in\mathcal M_t:\exists s,\ \operatorname{valid}_{j}^{\mathrm{cache},[s]}(t)=1,
\ D_{j\mid W_j^{[s]},F}^{2}(t)\le(\tau_F^{\mathrm{pr}}(t))^2\},
\] \[
\mathcal P_{\mathrm{Schur}}(t)
:=
\{j\in\mathcal M_t:\exists s,\ \operatorname{valid}_{j}^{\mathrm{cache},[s]}(t)=1,
\ L_j^{\mathrm{Schur}}(W_j^{[s]};t)\le\tau_{\mathrm{Schur}}^{\mathrm{pr}}(t)\}.
\] The uncapped trial universe is \[
\boxed{
\mathcal U_{\mathrm{trial}}(t)
:=
\mathcal N_1(t)
\cup
\mathcal C_{\mathrm{flat}}^{\mathrm{proj}}(t)
\cup
\mathcal C_{\mathrm{curv}}^{\mathrm{proj}}(t)
\cup
\mathcal P_F(t)
\cup
\mathcal P_{\mathrm{Schur}}(t).
}
\] The capped exact-trial set is \[
\boxed{
\mathcal C_{\mathrm{trial}}(t)
:=
\operatorname{Top}_{Q_{\mathrm{trial}}(t)}
\left(\preceq_{\mathrm{pr}};\ j\in\mathcal U_{\mathrm{trial}}(t)\right).
}
\] A conservative ordering is lexicographic in \[
\left(
\chi_{\mathrm{curv}}(j),
\ s_{\mathrm{geo}}(j;t),
\ L_j^{\mathrm{Schur}}(W_j^{[s_{\mathrm{geo}}]};t),
\ D_{j\mid W_j^{[s_{\mathrm{geo}}]},F}^{2}(t),
\ R_j^{\mathrm{cheap}}(t)
\right),
\] where
\(\chi_{\mathrm{curv}}(j):=\mathbf 1[s_{\mathrm{geo}}(j;t)\ge3]\), and
candidates with \(s_{\mathrm{geo}}=\infty\) are ranked by the cheap
order unless a scheduled compression pass overrides this. In all cases,
\(\mathcal C_{\mathrm{trial}}\) only selects what to test. It is not a
prune decision.

Optional cache-maintenance details such as block-prune Schur formulae,
monotonicity diagnostics, damped BFGS updates, and metric-coordinate
cache transport are auxiliary. They may improve ranking and shot
allocation, but they are not part of the deletion certificate.

\paragraph{Prune 3: exact remove--refit
ladder}\label{prune-3-exact-removerefit-ladder}

Prune 3 is the first stage that actually changes the trial scaffold. For
a candidate \(j\in\mathcal C_{\mathrm{trial}}(t)\), the controller
deletes \(j\), chooses a typed compensator window \(W_j^{[s]}\), and
refits only that surviving window while all other coordinates remain
fixed. The rung \(s\) controls how much compensating freedom is opened:
\(s=0\) is frozen deletion, \(s=1,2\) are flat-sector or correlated
recovery attempts, \(s=3\) opens noncommuting compensators, and \(s=4\)
is the broad terminal or scheduled-compression attempt.

The exact remove-refit ladder is \[
\boxed{
\delta E_j^{[s]}(t)
:=
\min_{\xi\in\Theta_{\mathcal O_{t,-j}^+}}
\left\{
E(\mathcal O_{t,-j}^+,\xi):
\xi_i=((\theta_t^+)_{-j})_i
\ \forall i\notin W_j^{[s]}(t)
\right\}
-
E(\mathcal O_t^+,\theta_t^+),
\qquad s\in\{0,1,2,3,4\}.
}
\] The first rung is frozen ablation, \[
\delta E_j^{[0]}(t)
=
E(\mathcal O_{t,-j}^+,(\theta_t^+)_{-j})
-
E(\mathcal O_t^+,\theta_t^+),
\] and, assuming comparable optimization at each rung, \[
\boxed{
\delta E_j^{[0]}(t)
\ge
\delta E_j^{[1]}(t)
\ge
\delta E_j^{[2]}(t)
\ge
\delta E_j^{[3]}(t)
\ge
\delta E_j^{[4]}(t)
\ge
\delta E_j^\star(t).
}
\] Append and prune use separate refit-breadth schedules. Append may
shrink its inherited old-coordinate fraction with maturity, while prune
broadens typed compensating windows with maturity: \[
\rho_W^{\mathrm{app}}(t)\downarrow,
\qquad
\rho_W^{\mathrm{pr}}(t)\uparrow.
\] In a shot-estimated route, typed rung \(s\) with shot budget \(N\)
uses the paired remove-refit estimator \[
\widehat\theta_{-j}^{[s,N]}
:=
\arg\min_{\xi\in\Theta_{\mathcal O_{t,-j}^+}}
\left\{
\widehat E_N(\mathcal O_{t,-j}^+,\xi):
\xi_i=((\theta_t^+)_{-j})_i
\ \forall i\notin W_j^{[s]}(t)
\right\},
\] \[
\boxed{
\widehat{\delta E}_{j}^{[s,N]}
:=
\widehat E_N(\mathcal O_{t,-j}^+,\widehat\theta_{-j}^{[s,N]})
-
\widehat E_N(\mathcal O_t^+,\theta_t^+).
}
\] Use paired Pauli-group measurement, shared random seeds, or
common-shot bookkeeping where possible, so that the variance of the
difference is estimated directly: \[
\widehat\sigma_{j}^{2}(s,N)
:=
\widehat{\operatorname{Var}}
\left[
\widehat E_N(\mathcal O_{t,-j}^+,\widehat\theta_{-j}^{[s,N]})
-
\widehat E_N(\mathcal O_t^+,\theta_t^+)
\right].
\] The upper and lower confidence bounds are \[
U_j^{[s,N]}
:=
\widehat{\delta E}_{j}^{[s,N]}
+
z_\alpha\widehat\sigma_j(s,N),
\qquad
L_j^{[s,N]}
:=
\widehat{\delta E}_{j}^{[s,N]}
-
z_\alpha\widehat\sigma_j(s,N).
\] These confidence bounds now pass to the final confirmation stage.

\paragraph{Prune 4: shot confirmation, recovery class, and
rollback}\label{prune-4-shot-confirmation-recovery-class-and-rollback}

Prune 4 converts the remove-refit trial into a committed deletion or an
exact rollback. It uses confidence bounds on the paired energy
regression, a retained admitted-gain guard, and an extra curvature guard
for any rung that opened noncommuting compensators. Cached projection,
amplitude collapse, and Schur proxies do not appear as deletion
certificates here; they have already done their job by choosing which
trial to spend budget on.

The recovery tolerance is absolute, shot-aware, and scale-aware: \[
\boxed{
\Delta_{\mathrm{pr}}^{\mathrm{tol}}(t)
:=
\max\!\left(
\Delta_{\mathrm{num}},
 c_{\mathrm{shot}}\sigma_E(t),
 c_{\mathrm{scr}}\Delta_{\mathrm{scr}}(t),
 \Delta_{\mathrm{chem}},
 c_{\mathrm{rel}}|E_t^+-E_{\mathrm{target}}|
\right),
}
\] dropping unavailable terms from the maximum. The retained-gain guard
is \[
\Gamma_{\mathrm{ret}}(j,t;s,N)
:=
\mathbf 1[\Delta_{\mathrm{adm}}(t)<\Delta_{\mathrm{ret,on}}]
\vee
\mathbf 1\!\left[
E_t^-
-
\widehat E_N(\mathcal O_{t,-j}^+,\widehat\theta_{-j}^{[s,N]})
\ge
\eta_{\mathrm{ret}}\Delta_{\mathrm{adm}}(t)
+
z_\alpha\widehat\sigma_j(s,N)
\right].
\] Define the stronger retained-gain guard
\(\Gamma_{\mathrm{ret}}^{\mathrm{strong}}(j,t;s,N)\) by replacing
\(\eta_{\mathrm{ret}}\) with
\(\eta_{\mathrm{ret}}^{\mathrm{strong}}\ge\eta_{\mathrm{ret}}\). Let
\(\operatorname{CompressionMode}_t\in\{0,1\}\) denote a scheduled
high-confidence compression pass. The curvature-compensated guard is \[
\boxed{
\begin{aligned}
\Gamma_{\mathrm{curv}}(j,t;s,N)
:=
&\ \mathbf 1[s\le2]
\vee
\mathbf 1[N\ge N_{\mathrm{hi}}]
\vee
\mathbf 1[s=4]
\vee
\mathbf 1\!\left[\frac{|W_j^{[s]}(t)|}{|I_{-j}|+\varepsilon}\ge\rho_{\mathrm{near}}\right]\\
&\vee
\mathbf 1\!\left[U_j^{[s,N]}\le \gamma_{\mathrm{curv}}\Delta_{\mathrm{pr}}^{\mathrm{tol}}(t)\right]
\vee
\mathbf 1\!\left[\Gamma_{\mathrm{ret}}^{\mathrm{strong}}(j,t;s,N)=1\right]
\vee
\mathbf 1[\operatorname{CompressionMode}_t=1],
\end{aligned}
}
\] where \(0<\gamma_{\mathrm{curv}}<1\). Any accepted rung that has
opened noncommuting compensators must be backed by at least one stronger
confirmation condition. Recovery through noncommuting compensators is
interpreted as curvature-mediated reparameterization, not proof of flat
redundancy.

The rung decision is \[
\boxed{
\mathsf D_j^{[s,N]}
:=
\begin{cases}
\mathrm{accept},
&
U_j^{[s,N]}
\le
\Delta_{\mathrm{pr}}^{\mathrm{tol}}(t)
\ \wedge\
\Gamma_{\mathrm{ret}}(j,t;s,N)=1
\ \wedge\
\Gamma_{\mathrm{curv}}(j,t;s,N)=1,
\\[1mm]
\mathrm{reject},
&
s=4,\ N=N_{\mathrm{hi}},\
L_j^{[s,N]}
>
\Delta_{\mathrm{pr}}^{\mathrm{tol}}(t),
\\[1mm]
\mathrm{escalate},
&
\text{otherwise.}
\end{cases}
}
\] If the first accepted rung has \(s\le2\), set
\(\operatorname{recovery\_class}_j=\mathrm{flat\_redundant}\). If the
first accepted rung has \(s\ge3\), set
\(\operatorname{recovery\_class}_j=\mathrm{curvature\_compensated}\). If
all scheduled rungs reject or defer without acceptance, set
\(\operatorname{recovery\_class}_j=\mathrm{failed}\). Escalation may
increase \(N\), move from \(W_j^{[s]}\) to \(W_j^{[s+1]}\), or defer to
a scheduled compression pass.

The final prune rule is \[
\boxed{
\operatorname{Prune}_j(t)=1
\Longleftrightarrow
\Gamma_{\mathrm{prune}}^{\mathrm{perm}}(t)=1
\ \wedge\
j\in\mathcal C_{\mathrm{trial}}(t)
\ \wedge\
\exists(s,N):\mathsf D_j^{[s,N]}=\mathrm{accept}.
}
\] Before remove-refit trials, store the rollback snapshot \[
\Xi_t^{\mathrm{pr}}
:=
(\mathcal O_t^+,\theta_t^+,E_t^+,\{c_i(t)\}_i,\text{current local optimizer state}).
\] If \(\Gamma_{\mathrm{prune}}^{\mathrm{perm}}(t)=0\) or
\(\mathcal C_{\mathrm{trial}}(t)=\varnothing\), no prune trial is
launched and the committed branch state remains
\((\mathcal O_t^+,\theta_t^+,E_t^+)\). Otherwise test candidates in
\(\preceq_{\mathrm{pr}}\) order up to the configured probe cap. If a
prune is accepted, commit \[
(\mathcal O_{t,-j}^+,\widehat\theta_{-j}^{[s,N]},\widehat E_N(\mathcal O_{t,-j}^+,\widehat\theta_{-j}^{[s,N]}))
\] as the new branch state, write
\((j,s,N,\operatorname{recovery\_class}_j)\) to the branch prune-event
log, and then delete the live metadata carried by coordinate \(j\). If
it is rejected, restore the branch state exactly to
\(\Xi_t^{\mathrm{pr}}\), set \(c_j(t)\leftarrow c_{\mathrm{cool}}\), and
continue according to the configured candidate budget. Hence prune
failure never invalidates an already committed admission, and rollback
is exact.

\subsubsection{11.5.2 Beam-adapt over scaffold alternatives: branch
state, pruning, and effective
selector}\label{beam-adapt-over-scaffold-alternatives-branch-state-pruning-and-effective-selector}

In this section beam-adapt is \textbf{purely structural}: each branch
carries an alternative ADAPT scaffold together with its locally refit
amplitudes. There is no time checkpoint, probe rollout, or
projected-dynamics integral in this section; those belong to the
separate time-dynamics beam surface developed later.

The cleanest way to read a live scaffold branch is \[
\mathfrak b=
\begin{aligned}[t]
\bigl(&\text{identity/history};\ \text{current scaffold and amplitudes};\\
&\text{current energy/depth};\ \text{cumulative selector totals/controller telemetry/status}\bigr).
\end{aligned}
\] More explicitly, write \[
\mathfrak b=\bigl(\operatorname{id}_b,\pi_b,\mathcal H_b,\mathcal O_b,\theta_b,E_b,d_b,S_b^{\mathrm{cum}},K_b^{\mathrm{cum}},\mathcal T_b^{\mathrm{ctrl}},\mathrm{status}_b,\mathrm{term}_b\bigr),
\] with \(\mathcal H_b=(r_{b,1},\dots,r_{b,d_b})\) and
\(r_{b,j}=(m_{b,j}^{\mathrm{adm}},p_{b,j}^{\mathrm{adm}})\). Here
\(\operatorname{id}_b\) and \(\pi_b\) are the branch and parent
identifiers, \(\mathcal H_b\) is the ordered history of admitted
records, \(\mathcal O_b\) is the scaffold currently carried by branch
\(b\), \(\theta_b\) is the refit parameter vector on that scaffold,
\(E_b=E(\theta_b;\mathcal O_b)\) is the current variational energy,
\(d_b\) is the branch ADAPT depth,
\(S_b^{\mathrm{cum}}=\sum_{j=1}^{d_b}s_{b,j}\) is the cumulative
admitted selector value,
\(K_b^{\mathrm{cum}}=\sum_{j=1}^{d_b}K_{b,j}^{\mathrm{sel}}\) is the
cumulative admitted selector burden, \(\mathcal T_b^{\mathrm{ctrl}}\)
stores branch-local controller telemetry, and
\((\mathrm{status}_b,\mathrm{term}_b)\) record live/terminal status and
termination label. The increment \(K_{b,j}^{\mathrm{sel}}\) is the
selector-cost contribution attached to the \(j\)-th admission on branch
\(b\); for a newly materialized child this contribution is the
branch-local Phase-3 cost \(K_3(r;\mathfrak b)\). Thus
\(K_b^{\mathrm{cum}}\) is cumulative selector-side cost bookkeeping
rather than a raw hardware count.

If beam search is entered from an incumbent scaffold
\((\mathcal O^{(0)},\theta^{(0)})\), the root branch is
\(\mathfrak b_{\mathrm{root}}=\bigl(\operatorname{id}_0,\varnothing,\mathcal H^{(0)},\mathcal O^{(0)},\theta^{(0)},E(\theta^{(0)};\mathcal O^{(0)}),d_0,S_{\mathrm{root}}^{\mathrm{cum}},K_{\mathrm{root}}^{\mathrm{cum}},\mathcal T_{\mathrm{root}}^{\mathrm{ctrl}},\mathrm{frontier},\varnothing\bigr)\),
with \(d_0=|\mathcal H^{(0)}|\),
\(S_{\mathrm{root}}^{\mathrm{cum}}=\sum_{j=1}^{d_0}s_j^{(0)}\), and
\(K_{\mathrm{root}}^{\mathrm{cum}}=\sum_{j=1}^{d_0}K_j^{\mathrm{sel},(0)}\),
where \(K_j^{\mathrm{sel},(0)}\) denotes the selector-cost contribution
attached to the \(j\)-th historical admission on the incumbent scaffold.
If the beam is started from a fresh scaffold, then
\(\mathcal H^{(0)}=\varnothing\), \(d_0=0\), and both cumulative sums
vanish.

For clock and budget, set \(t_b:=d_b\) and
\(D_{\mathrm{left}}(\mathfrak b):=d_{\max}-d_b\). Now apply the branch
lift rule: for any live-state quantity \(X\) from Phase 3 whose
definition is written in terms of
\((\mathcal H,\mathcal O,\theta,t,\mathcal T^{\mathrm{ctrl}})\), define
\(X(\mathfrak b):=X\big|_{(\mathcal H,\mathcal O,\theta,t,\mathcal T^{\mathrm{ctrl}})\mapsto(\mathcal H_b,\mathcal O_b,\theta_b,t_b,\mathcal T_b^{\mathrm{ctrl}})}\);
for any record-local quantity \(X(r;t)\), define
\(X(r;\mathfrak b):=X(r;t)\big|_{(\mathcal H,\mathcal O,\theta,t,\mathcal T^{\mathrm{ctrl}})\mapsto(\mathcal H_b,\mathcal O_b,\theta_b,t_b,\mathcal T_b^{\mathrm{ctrl}})}\).
In particular, \(R(\mathfrak b)\), \(u_{\mathrm{sat}}(\mathfrak b)\),
\(\widehat N_{\mathrm{rem}}(\mathfrak b)\), \(K_3(r;\mathfrak b)\), and
\(g_{\mathrm{lcb}}(r;\mathfrak b)\) are branch-lifted evaluations of
their previously defined live-state counterparts.

The Phase-3 chain is branch-local by construction: all of its
dependencies factor through the state tuple
\((\mathcal O_b,\theta_b,t_b,\mathcal T_b^{\mathrm{ctrl}})\). Hence each
live branch induces its own raw candidate-position surface
\(\mathcal R_b^{\mathrm{raw}}\), its own branch-local shortlist
\(\mathcal S_3(\mathfrak b)\), and its own retained admission set
\(\mathcal A_b\).

For stage and enabled family, let
\(\eta_b\in\{\mathrm{core/seed},\mathrm{residual}\}\) and \[
\mathcal G_{\eta_b}^{(3)}(\mathfrak b)=
\begin{cases}
\mathcal G_{\mathrm{core/seed}}^{(3)}(\mathfrak b),&\eta_b=\mathrm{core/seed},\\
\mathcal G_{\mathrm{residual}}^{(3)}(\mathfrak b),&\eta_b=\mathrm{residual}.
\end{cases}
\] Let
\(\mathcal P_{\mathrm{residual}}(\mathfrak b)\subseteq\mathcal G^{(3)}\)
denote the residual-only generator family on branch \(b\). Then the
controller-stage gate and branch-local available Phase-3 generator
family are \[
\Gamma_{\mathrm{stage},b}(m)=
\begin{cases}
1,&\eta_b=\mathrm{residual},\\
1,&\eta_b=\mathrm{core/seed}\ \text{and}\ m\notin\mathcal P_{\mathrm{residual}}(\mathfrak b),\\
0,&\eta_b=\mathrm{core/seed}\ \text{and}\ m\in\mathcal P_{\mathrm{residual}}(\mathfrak b),
\end{cases}
\qquad
\mathcal G_{\mathrm{avail}}^{(3)}(\mathfrak b)=\{m\in\mathcal G^{(3)}:\Gamma_{\mathrm{stage},b}(m)=1\}.
\] For scaffold length, append slot, and active probe window, let
\(n_b=|\mathcal O_b|\), \(p_{\mathrm{app}}(\mathfrak b)=n_b\), and
\(W_{\mathrm{act}}(\mathfrak b)\subseteq\{0,\dots,n_b-1\}\). The ordered
probe list is
\(\mathcal L_{\mathrm{probe}}(\mathfrak b):=\bigl(p_{\mathrm{app}}(\mathfrak b),0,W_{\mathrm{act}}(\mathfrak b)\bigr)\),
and the distinct retained probe positions are
\(\mathcal P_{\mathrm{probe}}(\mathfrak b)=\operatorname{Head}_{M_{\mathrm{probe}}}\!\Bigl(\operatorname{Dedup}_{\mathrm{ord}}\bigl(\mathcal L_{\mathrm{probe}}(\mathfrak b)\bigr)\Bigr)\subseteq\{0,\dots,n_b\}\).
For symmetry and admissible positions, let \[
\Gamma_{\mathrm{sym},b}(m,p)=
\begin{cases}
0,&\text{the symmetry audit for }(m,p;\mathcal O_b)\text{ returns a hard violation},\\
1,&\text{otherwise},
\end{cases}
\] so the branch-local raw candidate-position surface is \[
\mathcal R_b^{\mathrm{raw}}
:=
\{\,r=(m,p):m\in\mathcal G_{\mathrm{avail}}^{(3)}(\mathfrak b),\ p\in\mathcal P_{\mathrm{probe}}(\mathfrak b),\ \Gamma_{\mathrm{sym},b}(m,p)=1\,\}.
\] The resulting local shortlist and retained admission set are \[
\mathcal S_3(\mathfrak b)
:=
\operatorname{Top}_{N_3(t_b)}\!\left(S_3(\cdot;\mathfrak b);\mathcal R_b^{\mathrm{raw}}\right),
\qquad
\mathcal A_b
:=
\operatorname{Top}_{K_b}\!\left(S_3(\cdot;\mathfrak b);\{r\in\mathcal S_3(\mathfrak b):S_3(r;\mathfrak b)>0\}\right).
\] When a branch-local rerank shell is judged tie-degenerate, the
retained admission set may be replaced by a transient tie-shell
selector.\footnote{For example, if $r_{b,\max}\in\arg\max_{r\in\mathcal S_3(\mathfrak b)}S_3(r;\mathfrak b)$, one may define $\mathcal T_{\mathrm{tie}}(\mathfrak b):=\{\,r\in\mathcal S_3(\mathfrak b): S_3(r;\mathfrak b)\ge \max(\eta_{\mathrm{tie}}\,S_3(r_{b,\max};\mathfrak b),\,S_3(r_{b,\max};\mathfrak b)-\delta_{\mathrm{tie}})\,\}$ and replace $\mathcal A_b$ by a retained subset $\mathcal A_b^{\mathrm{tie}}\subseteq \mathcal T_{\mathrm{tie}}(\mathfrak b)$ with $1\le |\mathcal A_b^{\mathrm{tie}}|\le B_{\mathrm{tie}}(t_b)$, chosen by descending $S_3(\cdot;\mathfrak b)$.}

\textbf{(ii) Materialize either stop or one new admission.} The proposal
family is \(\mathcal Q(b)=\{\mathrm{stop}\}\cup\mathcal A_b\). The
symbol \(\mathrm{stop}\) terminates the scaffold as-is, while each
\(r=(m,p)\in\mathcal A_b\) inserts one new generator into the branch and
refits. If \(\mathcal O_b=(m_{b,1},\dots,m_{b,n_b})\) and
\(0\le p\le n_b\), then the scaffold insertion map is
\(\operatorname{Insert}(\mathcal O_b,m,p)=(m_{b,1},\dots,m_{b,p},m,m_{b,p+1},\dots,m_{b,n_b})\),
the insertion index transport is
\(I_p(j)=\begin{cases} j,&1\le j\le p,\\ j+1,&p<j\le n_b, \end{cases}\),
and the zero-lifted parameter vector is
\(\theta_b^{\uparrow(r)}=(\theta_{b,1},\dots,\theta_{b,p},0,\theta_{b,p+1},\dots,\theta_{b,n_b})\).

For contiguous inserted-window cap \(\omega_{\mathrm{new}}\), let
\(\omega_{\mathrm{eff}}^{(p)}:=\max\{1,\min(\omega_{\mathrm{new}},n_b+1)\}\)
and
\(\ell_p:=\max\!\left\{1,\min\!\left(p+1-\left\lfloor\frac{\omega_{\mathrm{eff}}^{(p)}-1}{2}\right\rfloor,\;n_b+2-\omega_{\mathrm{eff}}^{(p)}\right)\right\}\).
Then
\(W_{\mathrm{new}}^{(p)}=\{\ell_p,\;\ell_p+1,\;\ldots,\;\ell_p+\omega_{\mathrm{eff}}^{(p)}-1\}\),
so \(|W_{\mathrm{new}}^{(p)}|=\omega_{\mathrm{eff}}^{(p)}\) and
\(p+1\in W_{\mathrm{new}}^{(p)}\). Let the retained old-coordinate
fraction shrink with maturity according to
\(\rho_W(t):=\rho_W^{\min}+(\rho_W^{\max}-\rho_W^{\min})e_t\), with
\(0\le \rho_W^{\min}\le \rho_W^{\max}\le 1\), and define the
old-coordinate complement by
\(C_{\mathrm{old}}^{(p)}:=\{1,\dots,n_b+1\}\setminus W_{\mathrm{new}}^{(p)}\),
so \(|C_{\mathrm{old}}^{(p)}|=n_b+1-\omega_{\mathrm{eff}}^{(p)}\). Then
the retained old-coordinate count is
\(n_{\mathrm{old}}(t,p):=\left\lceil \rho_W(t)\,|C_{\mathrm{old}}^{(p)}|\right\rceil=\left\lceil \rho_W(t)\,\bigl(n_b+1-\omega_{\mathrm{eff}}^{(p)}\bigr)\right\rceil\).
For each current scaffold coordinate \(j\in\{1,\dots,n_b\}\), let
\(a_{b,j}^{\mathrm{adm}}\) denote the admission step at which that
coordinate entered the current scaffold on branch \(b\), with smaller
values meaning older coordinates, and let the insertion-lifted
admission-order vector be
\(a_b^{\uparrow(r)}:=(a_{b,1}^{\mathrm{adm}},\dots,a_{b,p}^{\mathrm{adm}},+\infty,a_{b,p+1}^{\mathrm{adm}},\dots,a_{b,n_b}^{\mathrm{adm}})\),
so the newly inserted slot is never selected as an old coordinate. The
retained old-coordinate window is then
\(W_{\mathrm{age}}^{(p)}:=\{\,j\in C_{\mathrm{old}}^{(p)}:\#\{\,i\in C_{\mathrm{old}}^{(p)}:(a_b^{\uparrow(r)})_i<(a_b^{\uparrow(r)})_j\,\}<n_{\mathrm{old}}(t,p)\,\}\),
namely the set of the \(n_{\mathrm{old}}(t,p)\) oldest old coordinates
outside \(W_{\mathrm{new}}^{(p)}\), and hence the inherited refit window
on the inserted scaffold is
\(W_r=W(r;t):=W_{\mathrm{new}}^{(p)}\cup W_{\mathrm{age}}^{(p)}\).

Then, for \(r=(m,p)\) and
\(\mathcal O_b\oplus_p m:=\operatorname{Insert}(\mathcal O_b,m,p)\), the
branch-local refit map is
\(\operatorname{Refit}_{W_r}(\theta_b;r):=\arg\min_{\widetilde\theta\in\Theta_{\mathcal O_b\oplus_p m}}\left\{E(\mathcal O_b\oplus_p m,\widetilde\theta):\ \widetilde\theta_j=(\theta_b^{\uparrow(r)})_j\ \forall j\notin W_r\right\}\).

The transfer map is
\(\mathcal T(\mathfrak b,\mathrm{stop})=\mathfrak b'\) with
\(\mathcal H_{b'}=\mathcal H_b\), \(\mathcal O_{b'}=\mathcal O_b\),
\(\theta_{b'}=\theta_b\), \(E_{b'}=E_b\), \(d_{b'}=d_b\),
\(S_{b'}^{\mathrm{cum}}=S_b^{\mathrm{cum}}\),
\(K_{b'}^{\mathrm{cum}}=K_b^{\mathrm{cum}}\),
\(\mathrm{status}_{b'}=\mathrm{terminal}\), and
\(\mathrm{term}_{b'}=\begin{cases} \mathrm{empty},&\mathcal A_b=\varnothing,\\ \mathrm{stop},&\mathcal A_b\neq\varnothing. \end{cases}\)
For \(r=(m,p)\), \(\mathcal T(\mathfrak b,r)=\mathfrak b'\) with
\(\mathcal O_{b'}=\operatorname{Insert}(\mathcal O_b,m,p)\),
\(\theta_{b'}=\operatorname{Refit}_{W_r}(\theta_b;r)\),
\(\mathcal H_{b'}=(\mathcal H_b,r)\),
\(E_{b'}=E(\theta_{b'};\mathcal O_{b'})\), \(d_{b'}=d_b+1\),
\(S_{b'}^{\mathrm{cum}}=S_b^{\mathrm{cum}}+S_3(r;\mathfrak b)\),
\(K_{b'}^{\mathrm{cum}}=K_b^{\mathrm{cum}}+K_3(r;\mathfrak b)\),
\(\mathrm{status}_{b'}=\begin{cases} \mathrm{terminal},&d_{b'}\ge d_{\max}\ \text{or}\ \mathcal A_{b'}=\varnothing,\\ \mathrm{frontier},&d_{b'}<d_{\max}\ \text{and}\ \mathcal A_{b'}\neq\varnothing, \end{cases}\),
and
\(\mathrm{term}_{b'}=\begin{cases} \mathrm{depth\_cap},&d_{b'}\ge d_{\max},\\ \mathrm{empty},&d_{b'}<d_{\max}\ \text{and}\ \mathcal A_{b'}=\varnothing,\\ \varnothing,&d_{b'}<d_{\max}\ \text{and}\ \mathcal A_{b'}\neq\varnothing. \end{cases}\)
where \(\mathcal A_{b'}\) is recomputed from step (i) on the child
branch itself. So a non-stop child is exactly one more ADAPT admission
applied to the parent scaffold, followed by the inherited-window refit
already defined above, and its cumulative \(K\)-coordinate is updated by
the admitted branch-local Phase-3 selector burden
\(K_3(r;\mathfrak b)\). If \(\mathcal A_b=\varnothing\), then
\(\mathcal Q(b)=\{\mathrm{stop}\}\) and the branch contributes only a
terminal child.

Branch-local rollback law: before materializing \(r=(m,p)\), store the
branch snapshot
\(\Xi_b=(\mathcal O_b,\theta_b,E_b,\mathcal H_b,S_b^{\mathrm{cum}},K_b^{\mathrm{cum}},\text{optimizer memory})\).
If the post-refit energy violates \(E_{b'}>E_b+\epsilon_{\mathrm{rb}}\),
the parameter rollback mode restores only the zero-lifted pre-refit
parameters on the enlarged scaffold, while the structural rollback mode
rejects the admission and restores \(\Xi_b\) exactly. Thus insertion,
beam branching, batching, and pruning remain problem-generic, while an
outer controller-tuning layer may tune the insertion trigger, rollback
mode, and \(\epsilon_{\mathrm{rb}}\) as policy parameters rather than
Hamiltonian-specific exceptions.

\textbf{(iii) Classify, deduplicate, and prune the scaffold children.}
The frontier and terminal descendants of branch \(b\) are
\(\mathfrak C_b^{\mathrm{front}}=\{\mathfrak b' : \mathfrak b'=\mathcal T(\mathfrak b,r),\ r\in\mathcal A_b,\ \mathrm{status}_{b'}=\mathrm{frontier}\}\)
and
\(\mathfrak C_b^{\mathrm{term}}=\{\mathfrak b' : \mathfrak b'=\mathcal T(\mathfrak b,q),\ q\in\mathcal Q(b),\ \mathrm{status}_{b'}=\mathrm{terminal}\}\).
Two branches are treated as duplicates when they represent the same
scaffold and essentially the same refit point, for example under the
fingerprint
\(\operatorname{rnd}_{10}(x):=10^{-10}\operatorname{round}(10^{10}x)\),
\(\operatorname{round}_{10}(\theta_b):=\bigl(\operatorname{rnd}_{10}(\theta_{b,1}),\dots,\operatorname{rnd}_{10}(\theta_{b,n_b})\bigr)\),
and
\(\mathrm{fp}(\mathfrak b)=\bigl(d_b,\operatorname{labels}(\mathcal O_b),\operatorname{round}_{10}(\theta_b)\bigr)\).
Among near-equivalent branches, keep the one with best lexicographic
prune key
\(\pi_{\mathrm{prune}}(\mathfrak b)=\bigl(E_b,-S_b^{\mathrm{cum}},K_b^{\mathrm{cum}},|\mathcal O_b|,\operatorname{labels}(\mathcal O_b),\operatorname{round}_{10}(\theta_b),\operatorname{id}_b\bigr)\).

The beam operators are therefore
\(\operatorname{Dedup}(\mathcal X)=\{\text{best branch per fingerprint under }\pi_{\mathrm{prune}}\}\)
and
\(\operatorname{Prune}(\mathcal X;B)=\operatorname{Top}^{\min}_B\!\left(\pi_{\mathrm{prune}};\operatorname{Dedup}(\mathcal X)\right)\).
Hence the live and terminal pools update by
\(\mathfrak F_{c+1}=\operatorname{Prune}\!\left(\bigcup_{b\in\mathfrak F_c}\mathfrak C_b^{\mathrm{front}};B_{\mathrm{live}}\right)\)
and
\(\mathfrak T_{c+1}=\operatorname{Prune}\!\left(\mathfrak T_c\cup\bigcup_{b\in\mathfrak F_c}\mathfrak C_b^{\mathrm{term}};B_{\mathrm{term}}\right)\).

\textbf{(iv) Choose the final scaffold branch.} After \(C\)
scaffold-beam rounds, the surviving candidate set is
\(\mathfrak W_{\mathrm{final}}=\mathfrak F_C\cup\mathfrak T_C\), and the
effective beam selector is
\(\mathfrak b_*=\arg\min_{\mathfrak b\in\mathfrak W_{\mathrm{final}}}\pi_{\mathrm{prune}}(\mathfrak b)\).

A compact scaffold-beam algorithm block is therefore \[
\begin{aligned}
\text{A: }&
\mathfrak F_0=\{\mathfrak b_{\mathrm{root}}\},
\qquad
\mathfrak T_0=\varnothing,\\
\text{B: }&
\mathcal R_b^{\mathrm{raw}}
\to
\mathcal S_3(\mathfrak b)
\to
\mathcal A_b
\qquad(\mathfrak b\in\mathfrak F_c),\\
\text{C: }&
\mathcal Q(b)=\{\mathrm{stop}\}\cup\mathcal A_b,
\qquad
\mathfrak C_b=\{\mathcal T(\mathfrak b,q):q\in\mathcal Q(b)\},\\
\text{D: }&
\mathfrak C_b^{\mathrm{front}}
=
\{\mathfrak b'\in\mathfrak C_b:\mathrm{status}_{b'}=\mathrm{frontier}\},
\qquad
\mathfrak C_b^{\mathrm{term}}
=
\{\mathfrak b'\in\mathfrak C_b:\mathrm{status}_{b'}=\mathrm{terminal}\},\\
\text{E: }&
\mathfrak F_{c+1}
=
\operatorname{Prune}\!\left(\bigcup_{b\in\mathfrak F_c}\mathfrak C_b^{\mathrm{front}};B_{\mathrm{live}}\right),
\qquad
\mathfrak T_{c+1}
=
\operatorname{Prune}\!\left(\mathfrak T_c\cup\bigcup_{b\in\mathfrak F_c}\mathfrak C_b^{\mathrm{term}};B_{\mathrm{term}}\right),\\
\text{F: }&
\mathfrak W_{\mathrm{final}}=\mathfrak F_C\cup\mathfrak T_C,
\qquad
\mathfrak b_*=\arg\min_{\mathfrak b\in\mathfrak W_{\mathrm{final}}}\pi_{\mathrm{prune}}(\mathfrak b).
\end{aligned}
\]

\subsection{11.6 Active HH pool surface and selected scaffold}

Identify the nested active HH pool with the retained generator images by
\(\Omega_{\mathrm{HH}}^{(k)}:=\mathcal G^{(k)}\) for \(k\in\{1,2,3\}\).
The active HH adaptive surface is therefore carried by
\(\Omega_{\mathrm{HH}}^{(3)}\subseteq\Omega_{\mathrm{HH}}^{(2)}\subseteq\Omega_{\mathrm{HH}}^{(1)}\),
with
\(\mathcal G_{\mathrm{adapt}}^{(k)}=\{G_m\}_{m\in\Omega_{\mathrm{HH}}^{(k)}}\).

If beam continuation is used, set
\(\mathcal O_*:=\mathcal O_{\mathfrak b_*}\) and
\(\theta_*^{\mathrm{adapt}}:=\theta_{\mathfrak b_*}\); otherwise
\((\mathcal O_*,\theta_*^{\mathrm{adapt}})\) denotes the final admitted
scaffold and its refit amplitudes on the main branch. Write
\(n_*:=|\mathcal O_*|\) for the final scaffold length, and write the
selected ordered scaffold as \(\mathcal O_*=(G_1,\dots,G_{n_*})\). For
each \(j\in\{1,\dots,n_*\}\), write
\(G_j=\sum_{a=1}^{q_j} c_{ja}P_{ja}\), where \(q_j\) is the number of
active nonidentity Pauli terms in the \(j\)-th selected generator,
\(P_{ja}\) is the \(a\)-th such Pauli term, and \(c_{ja}\in\mathbb R\)
its coefficient. Decompose the selected amplitudes as
\(\theta_*^{\mathrm{adapt}}=(\theta_{*,1}^{\mathrm{adapt}},\dots,\theta_{*,n_*}^{\mathrm{adapt}})\)
with \(\theta_{*,j}^{\mathrm{adapt}}\in\mathbb R^{q_j}\). Thus each
selected generator carries a parameter block rather than a singleton
scalar.\footnote{When block-valued selected generators are used, the clean index convention is $j$ for the scaffold block, $a$ for the coordinate inside block $j$, and $\mu=(j,a)$ for the flattened global coordinate. Metric, force, regularized normal, and tangent matrices are indexed by global coordinates $\mu,\nu$, while append and prune decisions may still operate on block labels $j$.}
The resulting state is
\(|\psi_*\rangle=U(\theta_*^{\mathrm{adapt}};\mathcal O_*)|\phi_0\rangle\),
so the active HH scaffold manifold selected by ADAPT may be written as
\(\mathcal M_{\mathrm{scaf}}(\mathcal O_*)=\{\,U(\vartheta;\mathcal O_*)|\phi_0\rangle:\vartheta\in\prod_{j=1}^{n_*}\mathbb R^{q_j}\,\}\).

\subsection{Appendix 11-S. Supplemental algebraic
telemetry}\label{appendix-11-s.-supplemental-algebraic-telemetry}

The canonical selector uses exact local relation counts to assign
algebraic lanes. The continuous quantities below are supplemental
telemetry only. They may be logged, used in lane-budget diagnostics,
used in prune-risk bookkeeping, or activated in ablation studies, but
they do not multiply \(S_{1,\mathrm{rank}}\), \(S_2\), or \(S_3\).

Let \(\varepsilon_{\mathrm{alg}}>0\) be the small denominator stabilizer
used in algebraic ratios. For a candidate-position record \(r=(m,p)\),
define \[
\begin{aligned}
\operatorname{local\_comm\_degree}(r)
&:=
\frac{
\sum_{n\in\mathcal N_{\mathrm{loc}}(m)}
\Omega_{mn}\mathsf C_{mn}
}{
\sum_{n\in\mathcal N_{\mathrm{loc}}(m)}
\Omega_{mn}+\varepsilon_{\mathrm{alg}}
},\\[1mm]
\operatorname{local\_noncomm\_degree}(r)
&:=
\frac{
\sum_{n\in\mathcal N_{\mathrm{loc}}(m)}
\Omega_{mn}(1-\mathsf C_{mn})
}{
\sum_{n\in\mathcal N_{\mathrm{loc}}(m)}
\Omega_{mn}+\varepsilon_{\mathrm{alg}}
},\\[1mm]
\operatorname{noncomm\_bridge\_score}(r)
&:=
\operatorname{clip}_{[0,1]}\!\left(
\frac{
\#\{(a,b):a\neq b,\ a,b\in\mathcal N_{\mathrm{loc}}(m),\
\mathsf C_{ab}=1,\ \mathsf C_{ma}=0,\ \mathsf C_{mb}=0\}
}{
\#\{(a,b):a\neq b,\ a,b\in\mathcal N_{\mathrm{loc}}(m),\ \mathsf C_{ab}=1\}
+\varepsilon_{\mathrm{alg}}
}
\right),\\[1mm]
\operatorname{flat\_cluster\_saturation}(r;t)
&:=
\sum_{j\in\mathcal O_t}
\mathbf 1[c(m_j)=c(m)]\,\mathbf 1[\Omega_{m_jm}=1].
\end{aligned}
\] The corresponding supplemental record fields are
\(\operatorname{comm\_cluster\_id}\),
\(\operatorname{local\_comm\_degree}\),
\(\operatorname{local\_noncomm\_degree}\),
\(\operatorname{noncomm\_bridge\_score}\), and
\(\operatorname{flat\_cluster\_saturation}\). These fields describe
algebraic context after the exact lane assignment has already been made.

For continuity with the retired continuous-prior route, let
\(\mathcal J_{\mathrm{rec}}(t)\) be the newest protected or active
scaffold coordinates considered by the selector, with recency weights
\(\omega_{j,t}\ge0\), and let
\(q_{0.9}^{\mathrm{cluster}}(t):=\operatorname{Quantile}_{0.9}\bigl(\{\operatorname{flat\_cluster\_saturation}(u;t):u\in\mathcal R_k(t)\}\bigr)\)
on the current phase-input record family. The retired telemetry
quantities are \[
\begin{aligned}
\rho_{\mathrm{nc}}^{\mathrm{rec}}(r;t)
&:=
\frac{
\sum_{j\in\mathcal J_{\mathrm{rec}}(t)}
\omega_{j,t}\,\Omega_{m_jm}\,(1-\mathsf C_{m_jm})
}{
\sum_{j\in\mathcal J_{\mathrm{rec}}(t)}
\omega_{j,t}\,\Omega_{m_jm}+\varepsilon_{\mathrm{alg}}
},\\[1mm]
\pi_{\mathrm{flat}}(r;t)
&:=
\operatorname{clip}_{[0,1]}\!\left(
\frac{\operatorname{flat\_cluster\_saturation}(r;t)}{
q_{0.9}^{\mathrm{cluster}}(t)+\varepsilon_{\mathrm{alg}}}
\right).
\end{aligned}
\] These ratios summarize recent noncommutation and flat-cluster
saturation as continuous diagnostics. They are not authoritative
selection terms.

\section{17A. Append--Prune McLachlan Dynamics for Driven Mixed
Fermion--Boson
Evolution}\label{a.-appendprune-mclachlan-dynamics-for-driven-mixed-fermionboson-evolution}

This chapter is the active mathematical source for the dynamics
manuscript. The method is Append--Prune McLachlan dynamics
(AP-McLachlan): a variational real-time evolution in which the ansatz
support is fixed on open time intervals and may be patched at discrete
time points by no edit, zero-initialized append, verified prune, or a
simultaneous append--prune exchange. Exact/reference trajectories are
diagnostic outputs only. Controller decisions use measurement-compatible
tangent geometry and observable estimates of the prepared variational
state.

The older hybrid-checkpoint, secant-forecast, turbulent-regime,
adaptive-integrator, and large-amplitude theta heuristics are not the
active route identity and should not guide Paper-II implementation. They
are archival derivation material only. The main law below intentionally
keeps time integration as a declared numerical setting; the paper-facing
mechanism is support-patch control through grouped McLachlan gain and
loss.

\subsection{17A.1 Time grid, ansatz support, and projected McLachlan
geometry}\label{a.1-time-grid-ansatz-support-and-projected-mclachlan-geometry}

Let the control time points be \[
0=  au_0<   au_1<\cdots<    au_T=T,
\qquad
 t_k=t_0+\tau_k,
\qquad
\Delta\tau_k=\tau_{k+1}-\tau_k .
\] At time point \(k\), the ordered variational support is
\(\mathcal S_k\). Its persistent block labels are \(b\in\mathcal B_k\).
A block may be a Pauli-polynomial generator \[
A_b=\sum_{\ell=1}^{m_b}c_{b\ell}P_{b\ell}.
\] The support is block-valued, while the McLachlan solve is
coordinate-valued. Chapter 17A therefore treats the coordinate
parameterization as part of the benchmark protocol. In the per
Pauli/polynomial-term convention, \[
I_{b,k}^{\mathrm{term}}:=\{(b,1),\ldots,(b,m_b)\},
\qquad
U_{k,\mathrm{term}}(\theta)
:=\prod_{b\in\mathcal B_k}^{\leftarrow}
\prod_{\ell=m_b}^{1}
\exp\!\left(-i\theta_{b\ell}c_{b\ell}P_{b\ell}\right).
\] In the per logical/macro-generator convention, \[
I_{b,k}^{\mathrm{log}}:=\{b\},
\qquad
U_{k,\mathrm{log}}(\alpha)
:=\prod_{b\in\mathcal B_k}^{\leftarrow}
\prod_{\ell=m_b}^{1}
\exp\!\left(-i\alpha_{b}c_{b\ell}P_{b\ell}\right),
\] with \(U_{k,\mathrm{log}}(\alpha)\) written as
\(\prod_b^{\leftarrow}\exp(-i\alpha_bA_b)\) only when the block is
represented as a grouped exact generator. The per-term tangent family is
indexed by \((b,\ell)\), whereas the logical tangent for block \(b\) is
the diagonal derivative through all component factors, \[
\bar T_{b,k}^{\mathrm{log}}
:=Q_{\psi_k}\partial_{\alpha_b}|\psi(\alpha;\mathcal S_k)\rangle
=Q_{\psi_k}\sum_{\ell=1}^{m_b}
\left.\partial_{\theta_{b\ell}}|\psi(\theta;\mathcal S_k)\rangle
\right|_{\theta_{b1}=\cdots=\theta_{bm_b}=\alpha_b}.
\] Thus the same block support can produce different tangent matrices,
different normal equations, and different active parameter counts.
Below, \(I_{b,k}\) and \(\theta_k\) denote the declared coordinate
convention, and the active coordinate support is \[
J_k:=\{1,\ldots,N_k\},
\qquad
N_k:=\sum_{b\in\mathcal B_k}|I_{b,k}| .
\] The prepared state is \[
|\psi_k\rangle
:=|\psi(\theta_k;\mathcal S_k)\rangle
=U_k(\theta_k)|\psi_{\mathrm{ref}}\rangle,
\qquad
U_k(\theta):=U(\theta;\mathcal S_k).
\] If AP-McLachlan is initialized from a completed static ansatz, the
handoff state is denoted \(|\psi_{\mathrm{initial}}\rangle\). A
same-seed dynamics row requires the declared coordinate vector to
reproduce that state up to ray tolerance: \[
\epsilon_{\mathrm{handoff}}
:=\min_{\varphi\in\mathbb R}
\left\|
U_k(\theta_0)|\psi_{\mathrm{ref}}\rangle
-e^{i\varphi}|\psi_{\mathrm{initial}}\rangle
\right\|_2
\le \tau_{\mathrm{handoff}}.
\] For the per logical/macro-generator convention this check uses the
logical seed coordinates \(\alpha_0\) themselves. It is not valid to
silently replace a per-term seed by averaged logical coordinates unless
the ray check above passes. During propagation the support labels are
fixed and only the driven Hamiltonian coefficients change: \[
H(t)=\sum_{\alpha\in\mathcal I_{\mathrm{stat}}}h_\alpha P_\alpha
+\sum_{\beta\in\mathcal I_{\mathrm{drv}}}c_\beta(t)P_\beta .
\] The drive profile is the scalar coefficient \(c_\beta(t)\); the drive
operator is the Pauli-polynomial channel multiplying it. If a driven
benchmark initializes the ansatz with a zero-angle drive-aligned block
\(D_\beta=\sum_\ell d_{\beta\ell}P_{\beta\ell}\), the Hamiltonian is
unchanged while the tangent plane gains \[
\left.\partial_{\phi_\beta}
e^{-i\phi_\beta D_\beta}|\psi_k\rangle
\right|_{\phi_\beta=0}
=-iD_\beta|\psi_k\rangle,
\qquad
\bar T_{\beta,k}^{\mathrm{drv}}
:=Q_{\psi_k}(-iD_\beta|\psi_k\rangle)
\] as an available McLachlan direction. This initialization changes
\(\mathcal S_0\) and the measured tangent geometry; it does not change
\(H(t)\). For a sampled time \(t_k\) or declared quadrature stage,
define the horizontal tangent block and the projected Schrodinger drift
\[
\bar T_k
:=Q_{\psi_k}\left[\partial_{\theta_{k,1}}|\psi_k\rangle,\ldots,
\partial_{\theta_{k,N_k}}|\psi_k\rangle\right],
\qquad
Q_{\psi_k}:=I-|\psi_k\rangle\langle\psi_k|,
\] \[
\bar b_k
:=-i\bigl(H(t_k)-\langle H(t_k)\rangle_{\psi_k}\bigr)|\psi_k\rangle .
\] The McLachlan metric, force, and regularized normal matrix are \[
G_k:=\Re(\bar T_k^\dagger\bar T_k),
\qquad
f_k:=\Re(\bar T_k^\dagger\bar b_k),
\qquad
K_k:=G_k+\Lambda_k .
\] The regularizer \(\Lambda_k\succeq0\) represents declared numerical
stabilization. It is not a new physical force and it is not a
route-selection score.

For any coordinate support \(S\) inside the active or zero-initialized
augmented tangent support, define the restricted metric, regularized
metric, and force \[
G_{k,SS}:=\Re(\bar T_{k,S}^{\dagger}\bar T_{k,S}),
\qquad
K_{k,SS}:=G_{k,SS}+\Lambda_{k,SS},
\qquad
f_{k,S}:=\Re(\bar T_{k,S}^{\dagger}\bar b_k),
\] and let the structural and regularized support velocities be \[
\dot\theta_{k,S}^{G}:=(G_{k,SS})_\tau^\dagger f_{k,S},
\qquad
\dot\theta_{k,S}^{K}:=(K_{k,SS})_{\tau,\varepsilon}^{\oplus} f_{k,S}.
\] Here \((\cdot)_\tau^\dagger\) is the supported geometric
pseudoinverse and \((\cdot)_{\tau,\varepsilon}^{\oplus}\) is the damped
supported inverse used by the implemented velocity. The scalar \[
f_{k,S}^{\top}(G_{k,SS})_\tau^\dagger f_{k,S}
=f_{k,S}^{\top}\dot\theta_{k,S}^{G}
=
(\dot\theta_{k,S}^{G})^\top G_{k,SS}\dot\theta_{k,S}^{G}
=\|\bar T_{k,S}\dot\theta_{k,S}^{G}\|_2^2
\] is the force vector overlapped with its metric-scaled response. The
retained normal equation is
\(G_{k,SS}\dot\theta_{k,S}^{G}=P_{k,S}^{G,\tau}f_{k,S}\), with
\(P_{k,S}^{G,\tau}:=G_{k,SS}(G_{k,SS})_\tau^\dagger\); if the force lies
inside the retained geometric range, this reduces to
\(f_{k,S}=G_{k,SS}\dot\theta_{k,S}^{G}\). The corresponding regularized
scalar \[
f_{k,S}^{\top}(K_{k,SS})_{\tau,\varepsilon}^{\oplus}f_{k,S}
=f_{k,S}^{\top}\dot\theta_{k,S}^{K}
\] is used consistently for propagation, append gain, prune loss, and
exchange scoring.

The propagated velocity on an open time interval is \[
\dot\theta_k
:=K_{k,\tau,\varepsilon}^{\oplus}f_k,
\qquad
K_k\dot\theta_k^{\mathrm{LS}}=f_k
\quad\text{for the formal least-squares solution.}
\] The parameter ODE between adjacent control time points is therefore
\[
\frac{d\theta}{dt}=F_{\mathcal S_k}(\theta,t),
\qquad
F_{\mathcal S}(\theta,t):=\bigl(G_{\mathcal S}(\theta)+\Lambda_{\mathcal S}(\theta,t)\bigr)_{\tau,\varepsilon}^{\oplus} f_{\mathcal S}(\theta,t),
\] with \(\mathcal S=\mathcal S_k\) held fixed until the next
support-patch decision. The benchmark may declare Euler, RK4, or another
fixed one-step method, but adaptive integrator switching is not part of
the active AP-McLachlan mechanism.

\subsection{17A.2 Structural residual and support
inadequacy}\label{a.2-structural-residual-and-support-inadequacy}

The structural miss is the component of the sampled horizontal
Schrodinger drift that is not represented by the current tangent support
before damping artifacts are interpreted as physics: \[
\begin{aligned}
\epsilon_{\mathrm{expr},k}^2
&:=\operatorname{clip}\!\left(
\|\bar b_k\|_2^2
-f_{k,J_k}^{\top}(G_{k,J_kJ_k})_\tau^\dagger f_{k,J_k},
0,
\|\bar b_k\|_2^2
\right),\\
\rho_{\mathrm{expr},k}
&:=\operatorname{clip}\!\left(
\frac{\epsilon_{\mathrm{expr},k}^2}{\|\bar b_k\|_2^2+\varepsilon},
0,1
\right).
\end{aligned}
\] The realized regularized residual is \[
\begin{aligned}
\epsilon_{\mathrm{real},k}^2
&:=\|\bar T_k K_{k,\tau,\varepsilon}^{\oplus}f_k-\bar b_k\|_2^2,\\
\rho_{\mathrm{num},k}
&:=\left[
\frac{\epsilon_{\mathrm{real},k}^2}{\|\bar b_k\|_2^2+\varepsilon}
-\rho_{\mathrm{expr},k}
\right]_+ .
\end{aligned}
\] Large \(\rho_{\mathrm{expr},k}\) means the ansatz tangent span lacks
physical drift directions and may justify append. Large
\(\rho_{\mathrm{num},k}\) means the solve is failing to realize
available tangent directions and should trigger solve repair, damping
review, support threshold review, or time-step reduction; it is not by
itself evidence for a new ansatz block.

A typical high-miss trigger is \[
\operatorname{HighMiss}_k
:=\mathbf 1\!\left[
\rho_{\mathrm{expr},k}>\tau_{\mathrm{miss}}
\ \wedge\
\epsilon_{\mathrm{expr},k}^2>\epsilon_{\mathrm{abs}}^{\mathrm{miss}}
\right]
\] with optional persistence over recent time points. Persistence may
reduce noise sensitivity, but it cannot import exact target-state
information into the decision rule.

\subsection{17A.3 Zero-initialized append and grouped
gain}\label{a.3-zero-initialized-append-and-grouped-gain}

An append candidate is a support record \(r=(m,p)\): generator block
\(m\) inserted at ordered position \(p\). If the block contains \(m_r\)
Pauli-polynomial components, then \(|I_{r,k}^+|=m_r\) in the per
Pauli/polynomial-term convention and \(|I_{r,k}^+|=1\) in the per
logical/macro-generator convention. Denote the new coordinate set by
\(I_{r,k}^+\). The append is zero-initialized, so \[
|\psi(\theta_k,0;\mathcal S_k^+(r))\rangle=|\psi(\theta_k;\mathcal S_k)\rangle,
\] and the state is continuous at the time point. The new block changes
the tangent plane available for subsequent McLachlan propagation.

The canonical normalized append gain is the grouped forward
extra-sum-of-squares gain in the regularized McLachlan objective: \[
\boxed{
\rho_{r,k}^{\mathrm{gain}}
:=
\frac{
\left[
f_{k,J_k\cup I_{r,k}^+}^{\top}
(K_{k,(J_k\cup I_{r,k}^+)(J_k\cup I_{r,k}^+)})_{\tau,\varepsilon}^{\oplus}
f_{k,J_k\cup I_{r,k}^+}
-
f_{k,J_k}^{\top}
(K_{k,J_kJ_k})_{\tau,\varepsilon}^{\oplus}
f_{k,J_k}
\right]_+
}{
\|\bar b_k\|_2^2+\varepsilon
} .
}
\] For a batched append decision \(R_k^+\), define \[
I_{R,k}^+:=\bigcup_{r\in R_k^+}I_{r,k}^+,
\qquad
J_{R,k}^+:=J_k\cup I_{R,k}^+ .
\] The joint batched gain is \[
\boxed{
\rho_{R,k}^{\mathrm{gain}}
:=
\frac{
\left[
f_{k,J_{R,k}^+}^{\top}
(K_{k,J_{R,k}^+J_{R,k}^+})_{\tau,\varepsilon}^{\oplus}
f_{k,J_{R,k}^+}
-
f_{k,J_k}^{\top}
(K_{k,J_kJ_k})_{\tau,\varepsilon}^{\oplus}
f_{k,J_k}
\right]_+
}{
\|\bar b_k\|_2^2+\varepsilon
} .
}
\] The singleton formula is recovered by \(R_k^+=\{r\}\). Both
quantities measure the amount of regularized horizontal drift newly
explained after the old coordinates and candidate coordinates are
considered together under the same retained-support, ridge, and damping
convention. Static hardware cost may rank candidate evaluation, \[
\Xi_{r,k}^{\mathrm{app}}
:=\frac{\rho_{r,k}^{\mathrm{gain}}}{(\kappa_{r,k}^{\mathrm{hw}})^{\alpha_{\mathrm{app}}}},
\] but the cost-pressure score is not the admission quantity. A
sub-threshold tangent gain is not promoted by being cheap.

For a candidate tangent block \(\bar U_{r,k}\), set \[
B_{r,k}:=\Re(\bar T_k^\dagger\bar U_{r,k}),
\qquad
C_{r,k}:=\Re(\bar U_{r,k}^\dagger\bar U_{r,k}),
\qquad
q_{r,k}:=\Re(\bar U_{r,k}^\dagger\bar b_k).
\] Under the same old-support inverse, ridge, damping, and retained
spectral support as the propagated velocity, the grouped gain can be
evaluated through the Schur block \[
w_{r,k}:=q_{r,k}-B_{r,k}^\top K_{k,\tau,\varepsilon}^{\oplus}f_k,
\qquad
S_{r,k}:=C_{r,k}+\Lambda_r-B_{r,k}^\top K_{k,\tau,\varepsilon}^{\oplus}B_{r,k},
\] \[
\Delta_{r,k}^{\mathrm{Schur}}
:=w_{r,k}^\top(S_{r,k})_\tau^\dagger w_{r,k}.
\] The Schur expression is a compact evaluation or scout/confirm
diagnostic for the grouped gain. The direct grouped before--after
quantity \(\rho_{r,k}^{\mathrm{gain}}\) remains the definition of append
acceptance when whitening, thresholding, or compression makes the
shortcut non-identical.

An append is admitted only from present-time local data: \[
\operatorname{Admit}_{\mathrm{app}}(r,k)=1
\iff
\operatorname{HighMiss}_k=1
\ \wedge\
\rho_{r,k}^{\mathrm{gain}}\ge\tau_{\mathrm{gain}}
\ \wedge\
S_{r,k}^{\mathrm{conf}}\ge\tau_{\mathrm{conf}} .
\] The confirmation score may include Schur consistency, candidate
independence, direction-change control, and measurement budget, but it
may not use exact future fidelity, exact evolved states, or benchmark
reference trajectories.

\subsection{17A.4 Verified prune and grouped deletion
loss}\label{a.4-verified-prune-and-grouped-deletion-loss}

A prune proposal deletes a persistent active block \(b\) with coordinate
set \(I_{b,k}\). Deletion is the backward grouped companion of append.
The canonical normalized full-support deletion loss is \[
\boxed{
L_{b,k}^{\mathrm{full}}
:=
\frac{
\left[
f_{k,J_k}^{\top}
(K_{k,J_kJ_k})_{\tau,\varepsilon}^{\oplus}
f_{k,J_k}
-
f_{k,J_k\setminus I_{b,k}}^{\top}
(K_{k,(J_k\setminus I_{b,k})(J_k\setminus I_{b,k})})_{\tau,\varepsilon}^{\oplus}
f_{k,J_k\setminus I_{b,k}}
\right]_+
}{
\|\bar b_k\|_2^2+\varepsilon
} .
}
\] For a prune batch \(B_k^-\), define \[
I_{B,k}^-:=\bigcup_{b\in B_k^-}I_{b,k},
\qquad
J_{B,k}^-:=J_k\setminus I_{B,k}^- .
\] The batched full-support deletion loss is \[
\boxed{
L_{B,k}^{\mathrm{full}}
:=
\frac{
\left[
f_{k,J_k}^{\top}
(K_{k,J_kJ_k})_{\tau,\varepsilon}^{\oplus}
f_{k,J_k}
-
f_{k,J_{B,k}^-}^{\top}
(K_{k,J_{B,k}^-J_{B,k}^-})_{\tau,\varepsilon}^{\oplus}
f_{k,J_{B,k}^-}
\right]_+
}{
\|\bar b_k\|_2^2+\varepsilon
} .
}
\] The singleton formula is recovered by \(B_k^-=\{b\}\). The loss is
computed using the same support convention, damping, ridge, and
measurement-compatible tangent data as propagation and append. Small
coordinate amplitude, old age, pairwise commutation, or a cheap local
proxy can nominate a candidate, but none of those quantities certifies
deletion.

Prune acceptance requires a reduced-state projection rather than raw
coordinate dropping. Let \(\mathcal S_k^{(-b)}\) be the support with
block \(b\) removed. The reduced parameter point is \[
\theta_k^{(-b)}
:=\arg\min_{\xi\in\Theta_{\mathcal S_k^{(-b)}}}
\left[
 d_{\mathrm{FS}}^2\!\left(\psi(\xi;\mathcal S_k^{(-b)}),\psi(\theta_k;\mathcal S_k)\right)
 +\lambda_Y\|Y(\xi;\mathcal S_k^{(-b)})-Y_k\|_{\Sigma_Y^{-1}}^2
 +\lambda_v\|v(\xi;\mathcal S_k^{(-b)})-v_k\|_2^2
\right],
\] where \(Y_k\) denotes the tracked observable bundle and
\(v_k=\bar T_k\dot\theta_k\) is the represented horizontal velocity. In
strict measurement-compatible routes, every component of \(Y_k\) used
here is estimated from the incumbent or candidate prepared circuit
state; exact/reference observables remain diagnostic side channels.

The prune acceptance predicate has the form \[
\operatorname{Accept}_{\mathrm{pr}}(b,k)=1
\iff
\Gamma_{\mathrm{perm}}^{\mathrm{pr}}(b,k)=1
\ \wedge\
L_{b,k}^{\mathrm{full}}\le\tau_{\mathrm{loss}}^{\mathrm{pr}}
\ \wedge\
\Delta_{\mathrm{ray},b,k}^{\mathrm{pr}}\le\tau_{\mathrm{ray}}^{\mathrm{pr}}
\ \wedge\
\Delta\rho_{b,k}^{\mathrm{pr}}\le\tau_{\rho}^{\mathrm{pr}}
\ \wedge\
\mathfrak C_{b,k}^{\mathrm{shadow}}=1
\ \wedge\
\operatorname{Persist}_{b,k}^{\mathrm{pr}}=1 .
\] Here \(\Gamma_{\mathrm{perm}}^{\mathrm{pr}}\) enforces age, cooldown,
protected-block, solve-health, and candidate-budget rules. The ray and
differential-miss terms prevent state or tangent-span damage. The shadow
certificate propagates the incumbent and reduced candidate for a short
horizon under the same declared integrator and Hamiltonian sampling,
then checks observable no-harm: \[
\mathfrak C_{b,k}^{\mathrm{shadow}}
:=\mathbf 1\!\left[
\max_{1\le h\le H_{\mathrm{pr}}}
\|Y_{k+h}^{(-b)}-Y_{k+h}^{\mathrm{inc}}\|_{\Sigma_Y^{-1}}
+\lambda_{\Delta Y}
\max_{1\le h\le H_{\mathrm{pr}}}
\|\Delta Y_{k+h}^{(-b)}-\Delta Y_{k+h}^{\mathrm{inc}}\|_{\Sigma_{\Delta Y}^{-1}}
\le\tau_{\mathrm{shadow}}^{\mathrm{pr}}
\right].
\] Rejected prune trials roll back exactly and place the tested block on
cooldown. Accepted prune events commit the projected reduced state
\(\bigl(\mathcal S_k^{(-b)},\theta_k^{(-b)}\bigr)\).

\subsection{17A.5 Support-exchange patch and event
map}\label{a.5-support-exchange-patch-and-event-map}

The structural AP-McLachlan object at time point \(k\) is the support
patch \[
\nu_k=(B_k^-,R_k^+),
\] where \(B_k^-\) is a set of active blocks proposed for deletion and
\(R_k^+\) is a set of candidate blocks proposed for zero-initialized
insertion. Empty choices recover stay, pure append, and pure prune.
Nonempty choices on both sides give an exchange patch that replaces
stale tangent coordinates with new tangent coordinates in one event.

For a patch, define \[
I_{\nu,k}^-:=\bigcup_{b\in B_k^-}I_{b,k},
\qquad
I_{\nu,k}^+:=\bigcup_{r\in R_k^+}I_{r,k}^+,
\qquad
J_k^\nu:=(J_k\setminus I_{\nu,k}^-)\cup I_{\nu,k}^+,
\] \[
\Delta_{\nu,k}^{\mathrm{patch}}
:=
\frac{
f_{k,J_k^\nu}^{\top}
(K_{k,J_k^\nu J_k^\nu})_{\tau,\varepsilon}^{\oplus}
f_{k,J_k^\nu}
-
f_{k,J_k}^{\top}
(K_{k,J_kJ_k})_{\tau,\varepsilon}^{\oplus}
f_{k,J_k}
}{
\|\bar b_k\|_2^2+\varepsilon
} .
\] The separate append and prune criteria remain hard side constraints:
every appended block must pass grouped gain and confirmation, every
deleted block must pass full-support loss and nonlinear verification,
and the combined patch score is recomputed on \(J_k^\nu\). Exchange can
therefore be implemented conservatively after pure append and pure prune
are tested; if the combined support check disagrees with the side
scouts, the controller falls back to pure append, pure prune, or stay.

The structural map is \[
\mathsf S_{\nu_k}(\mathcal S_k,\theta_k)=
\begin{cases}
(\mathcal S_k,\theta_k),&B_k^-=\varnothing,\ R_k^+=\varnothing,\\
(\mathcal S_k^+(R_k^+),\theta_k\oplus 0),&B_k^-=\varnothing,
\ R_k^+\ne\varnothing,\\
(\mathcal S_k^{(-B_k^-)},\theta_k^{(-B_k^-)}),&B_k^-\ne\varnothing,
\ R_k^+=\varnothing,\\
((\mathcal S_k^{(-B_k^-)})^+(R_k^+),\theta_k^{(-B_k^-)}\oplus0),&B_k^-\ne\varnothing,
\ R_k^+\ne\varnothing .
\end{cases}
\] After the patch is committed, the next interval is propagated by the
declared fixed integrator: \[
(\mathcal S_k,\theta_k)
\xrightarrow{\mathrm{diagnose/select}}
\nu_k^\star
\xrightarrow{\mathsf S_{\nu_k^\star}}
(\mathcal S_k^\nu,\theta_k^{\nu,+})
\xrightarrow{\operatorname{Step}_{\mathrm{declared}}}
(\mathcal S_{k+1},\theta_{k+1}).
\] Solve repair has priority over structural edits when the tangent
solve is unsupported, badly conditioned, or dominated by numerical
damping. Exact benchmarking may be computed during or after the run for
plots, errors, spectra, and audits, but it has no code path into
\(\nu_k^\star\) in a strict AP-McLachlan row.

\subsection{17A.6 Paper-facing evidence
contract}\label{a.6-paper-facing-evidence-contract}

The dynamics manuscript should describe the active route as AP-McLachlan
support-patch control over time points or time iterations. Legacy
implementation fields may still contain \texttt{checkpoint} for
compatibility, but the mathematical object is the time-indexed support
patch \(\nu_k=(B_k^-,R_k^+)\) and the projective McLachlan velocity on
the active variational manifold.

Every same-seed AP-McLachlan or fixed-McLachlan row must record its
coordinate parameterization. The two supported labels are per
Pauli/polynomial term and per logical/macro generator; implementation
records may abbreviate these as \texttt{per\_pauli\_term} and
\texttt{logical\_shared}. The record should expose the active coordinate
count, the logical block count, and the underlying Pauli-component
count, because appending a block changes \(N_k\) by \(m_r\) in the
per-term convention and by one in the logical convention. Rows using
different coordinate parameterizations are different variational tangent
geometries unless the benchmark explicitly presents them as a
coordinate-mode ablation.

Main-text claims should stay inside the demonstrated evidence surface:

\begin{itemize}
\tightlist
\item
  append is supported when selected rows show accepted zero-initialized
  tangent-block additions and improved grouped McLachlan support score;
\item
  prune is supported only when selected rows show committed verified
  deletions under the full-support loss, projection, shadow,
  persistence, cooldown, and rollback checks;
\item
  drive-aligned ansatz initialization is reported as a zero-angle
  support augmentation and not as a Hamiltonian modification;
\item
  logical/macro-generator rows are supported only when the handoff
  parity check against \(|\psi_{\mathrm{initial}}\rangle\) passes;
  otherwise the benchmark must fail loudly or use the per-term
  coordinate convention;
\item
  integrator switching, secant forecasts, turbulent-regime labels, and
  large-amplitude theta heuristics are archival diagnostics, not active
  Paper-II implementation guidance;
\item
  exact/reference trajectories are valid diagnostic comparators for
  energy, observables, fidelity, spectra, and plotting, but are
  forbidden as controller-decision inputs in strict AP-McLachlan runs.
\end{itemize}

For the present same-seed dynamics surface, the paper-facing language
should reflect the measured support-edit pattern: append can be
discussed as an active AP-McLachlan mechanism; pruning should be
presented as a verified controller capability and ablation surface
unless the selected evidence rows actually commit prune events.

\end{document}
